% 04d-baselines.tex -- beats 7-9 and the ledger. sec:res-field has NO claim
% block, deliberately. Source: Investigation-v4.tex 2392-2938. Numbers verbatim
% from v1; two pending-audit flags. "the bar" = training-only drug_lookup;
% "the reference" = within-well split-half; "the threshold" = 0.80.

\beatsection{R}{7}{A lookup table wins}{sec:res-lookup}
\label{sec:methods-baselines}
\claimlede{The model reads the drug and recovers a small fraction of its
residual.}

\begin{question}[7]
What does a drug-response prediction have to beat before it is worth having? The
instrument is a ladder of predictors scored through one pipeline, the lowest rung
of which needs no model at all, only a table of what each drug did elsewhere.
That table wins.
\end{question}

\noindent
A model can respond to the drug in its prompt and still predict badly, so the
pre-registered bar was never ``beats chance''. Residual targets are plausibly
memorisable, and a predictor that looks up what this drug did elsewhere needs no
model.\gloss{lookup}{Predicts this drug's residual as its mean residual in the
\emph{other} cell lines.} Two objections to \cref{sec:res-blind}'s
drug-blindness verdict come first, and neither needs a checkpoint.

% ---------------------------------------------------------------------------
\paragraph{Is there anything to know?} A ladder of mean-shift baselines, each
predicting the average shift of a progressively narrower group, answers that.
It is scored in the expression frame on $\DEdr$, so its rows may not be set
beside any $\NIR$ row below. On tier 2 the global mean scores $0.056$; cell line
raises it about two-and-a-half-fold to $0.142$; mechanism does not move it
($0.052$); both together are no better than cell line alone ($0.120$ against
$0.142$). Tiers 1 and 3 agree. The informative conditioning variable here is the
cellular context, not the drug class, the same asymmetry \cref{sec:res-used}
finds inside the model's activations. Mechanism does carry signal once the
generic programme is subtracted (\cref{sec:res-scope}); in full-profile space the
generic response swamps it.

\paragraph{Does the encoding keep the drug?} A sentence keeps only the rank
order of genes. If the drug lived in the discarded magnitudes, this chapter's
diagnosis would be attacking the wrong thing. It does not. At single-cell
resolution, cosine, Pearson and Spearman similarities on true normalised
expression all sit within $0.04$ of chance
(\ci{0.529\textrm{--}0.531}{0.515}{0.543}), the same near-chance regime as rank.
What recovers the signal is aggregation. At fifteen cells the control-referenced
shift cosine reaches \ci{0.793}{0.775}{0.809}. So the re-encoding of
\cref{sec:res-encoding} cannot be read as restoring magnitudes; it keeps rank
throughout and changes only what is ranked. The same measurement also bounds
what a single cell can deliver. The single-cell same-drug versus different-drug
gap is only
\ci{+0.006}{+0.003}{+0.008}, so these instruments do not resolve drug identity
from one observed cell at this sequencing depth and $946$-gene panel.

% ---------------------------------------------------------------------------
\paragraph{What makes the ladder fair} Each baseline in \cref{tab:lookup}, set
out against chance and the reference in \cref{fig:ladder}, predicts the
drug-specific residual through one pipeline, on the same truths, encoding and
comparison set. Two disciplines keep it honest. The lookups are
fitted from training conditions only; fitted from the full cache, a lookup
scoring a held-out condition would read residuals the model was never shown, an
oracle and not a baseline, so that construction is excluded. And
\texttt{drug\_lookup\_1} averages a chosen cell line's conditions, because $62\%$
of drugs are multi-dose and a single condition would confound ``different cell
line'' with ``different dose''. \texttt{moa\_lookup} is not an oracle in any
split either. Its partner pool is restricted to \emph{training} conditions in the
same cell line, so a held-out condition never enters the prediction made for it.
The unrestricted-pool defect audited in \cref{sec:res-scope} was in the channel
gate; the restriction was already in force here, so the $n = \num{218}$ row is on
the restricted side of that repair.

\begin{defn}{Coverage of the reference}
A predictor's \emph{coverage} is the fraction of the distance from chance to the
within-well split-half precision reference that it closes,
\[
  (\NIR_{\text{arm}} - 0.50)\,/\,(\NIR_{\text{ref}} - 0.50).
\]
Coverage above $100\%$ therefore means the arm scores above the reference, a
statement about the reference and not a headroom estimate. Reference and coverage
are computed on the same support as the arm they describe.
\end{defn}

% ---------------------------------------------------------------------------
\begin{figure}[htbp]
% NOT \begin{fullwidth}. Laid out at the 142mm body block, which returns this
% page's margin notes: the MoA column carries two marks in the bottom sixth and
% the upper two-thirds of a 174mm strip would be blank, so the width is not
% earned. Body only; the tikzpicture lives in figs/ladder.tex, and every number
% it draws is listed with its source pointer in figs/fig-ladder.json.
\centering
% ===========================================================================
% figs/ladder.tex -- fig:ladder, reworked for thesis_v2.
% FIGURE BODY ONLY. No preamble, no document environment. \input it from
% inside a plain figure float (NOT fullwidth) in Sections/04d-baselines.tex.
%
% This is the SUPPORT figure: three columns at three different supports
% (n = 1394 / 1192 / 218), each showing where its arms sit between chance and
% the within-well split-half precision reference. It is deliberately distinct
% from fig:two-scales, which is the FRAME figure. The argument is unchanged.
%
% WHAT CHANGED FROM THE INLINE v2 VERSION (five spec'd fixes, plus one number
% correction that is NOT cosmetic -- see PROVENANCE note 3 below):
%
%   (1) The chance rule at 0.50 struck through the word "generic", whose label
%       sat at its own value of exactly 0.500. The caption states that labels a
%       rule would strike are lifted clear; "generic" had slipped through it.
%       Its label is now lifted BELOW the rule (0.500 -> 0.478) with a leader.
%   (2) The 115.7% brace and its label sat hard against the col-1/col-2
%       divider. The brace stem is now at x=6.60, which is 2.15cm (61pt) clear
%       of the divider at x=4.45, and its label's left edge is 0.98cm (28pt)
%       clear of it. The tightest thing to the divider is now the 13.9% label,
%       whose left edge is 0.475cm (13.5pt) clear -- the spec asked for >=6pt
%       and every part of the brace apparatus beats that. The braces are NESTED,
%       short 13.9% brace outboard at x=5.95, tall 115.7% brace inboard at
%       x=6.60, with each label at its own brace's height. That is the only
%       nesting in which neither label crosses the other's stem: the short
%       brace exists only over y in [0.500,0.549], so the tall brace's label,
%       sitting at y=0.71, has nothing in its way.
%   (3) The four left-column leaders all left the same 3-dot blob at
%       0.506/0.501/0.500 (0.8mm and 0.16mm apart on this scale, against a
%       1.0mm dot). control_copy / random / generic are now dodged sideways to
%       x = 0.30 / 0.58 / 0.86 so each leader leaves its own dot, and the four
%       label heights (0.645 / 0.583 / 0.532 / 0.478) fan the leaders to four
%       clearly different angles. Horizontal position inside a column carries
%       no meaning; the caption says so.
%   (4) \begin{fullwidth} DROPPED. Laid out at 142mm (the v2 body block), which
%       returns that page's margin notes. Drawn extent is 13.96cm of 14.20cm
%       (measured: 400.4pt of the 404.03pt block). Nothing is \resizebox'd:
%       1cm here is 1cm on the page, and 9pt type prints at 9pt.
%   (5) Label typography reduced from four registers to three, and the key is
%       stated in the caption:
%             \ttfamily ink      an arm, named as tab:lookup names it
%             roman     ink      prose ("within-well reference")
%             roman     accent   the model
%       v1/inline-v2 set "model" in \textbf + accent, a fourth register with no
%       key, and split the identifiers arbitrarily between ink and quietink.
%       quietink is now furniture only (ticks, support sizes, axis name).
%   Also, neither spec'd nor argument-bearing:
%     - the y axis had no name. It now carries one (NIR, RESID frame), set at
%       the top of the tick column so it costs the figure no width;
%     - gridlines are broken at the dividers. The caption forbids reading
%       across columns and an unbroken full-width rule invites exactly that;
%     - every dot gets a white halo, so the two arms that sit ON the chance
%       rule (random 0.501, generic 0.500) read as marks rather than as
%       thickenings of the rule. Same device as the interrupted zero rule in
%       figs/comparators.tex.
%
% PROVENANCE of every number drawn is in fig-ladder.json, sibling to this file.
% Source for all of it: RESULTS_cluster/re_v3.json.
%   1. Column 1 (n=1394) reads means/{ceiling,model,control_copy,random,generic}.
%   2. Column 2 (n=1192) reads means/common_support/{model,ceiling,drug_lookup}
%      and means/drug_lookup_1, and the two coverage figures read
%      means/common_support/{coverage_model,coverage_lookup}.
%   3. Column 3 (n=218). moa_lookup = means/moa_lookup = 0.5520. The model mark
%      in this column was 0.549 -- that is the model's COMMON-SUPPORT value
%      reused on a different support, and it is wrong. The model's mean over
%      the 218 MoA-support conditions is 0.5956 (mean of records[].model over
%      the 218 records that carry a moa_lookup value). Check: 0.595575 minus
%      0.552015 = 0.043560, which reproduces means/vs_baselines/moa_lookup/gap
%      = 0.0435594 to six decimals -- the same +0.0436 the tab:lookup caption
%      quotes. At 0.549 the figure put the model BELOW moa_lookup, the wrong
%      sign for a gap the text states as positive. The mark is drawn at 0.596.
%      No claim moves: the two-way interval [-0.0480,+0.1352] still spans zero,
%      so moa_lookup is still "not separated from the model".
% ===========================================================================
\begin{tikzpicture}[x=1cm, y=72mm, every node/.style={inner sep=0pt}]

  % --- value-to-position map. One map, one axis, all three columns. ---------
  \def\lo{0.47}\def\hi{0.95}
  \def\Y#1{{(#1-\lo)/(\hi-\lo)}}

  % --- the three registers of label type, and the furniture sizes. ---------
  % 9pt, not \scriptsize (8pt): this figure is drawn at its final size and is
  % never scaled, so its type may sit just under the 11pt caption rather than
  % being shrunk below it.
  \def\marklab{\fontsize{9}{10.5}\selectfont}
  \def\ticklab{\fontsize{8}{9}\selectfont}
  \def\headlab{\sffamily\fontsize{9}{10.5}\selectfont\bfseries}
  \def\subhead{\sffamily\fontsize{8}{9.5}\selectfont}

  % --- column geometry (cm): the two dividers and the right edge. ----------
  \def\Da{4.45}\def\Db{10.80}\def\R{13.32}

  % --- drawing macros. Order of use is leaders, then dots, then labels, so a
  %     dot and its halo always cover the tail of its own leader. -----------
  \newcommand{\lead}[5]{\draw[#5!45, line width=0.3pt]
     (#1,\Y{#2}) -- (#3-0.06,\Y{#4});}
  \newcommand{\dotat}[3]{\fill[white] (#1,\Y{#2}) circle (2.1pt);
     \fill[#3] (#1,\Y{#2}) circle (1.4pt);}
  \newcommand{\labat}[4]{\node[anchor=west, text=#3, font=\marklab]
     at (#1,\Y{#2}) {#4};}

  % --- gridlines, broken at the dividers ------------------------------------
  \foreach \v in {0.60,0.70,0.80,0.90}{
    \draw[rulegrey!45, line width=0.3pt] (-0.10,\Y{\v}) -- (4.32,\Y{\v});
    \draw[rulegrey!45, line width=0.3pt] ( 4.58,\Y{\v}) -- (10.67,\Y{\v});
    \draw[rulegrey!45, line width=0.3pt] (10.93,\Y{\v}) -- (\R,\Y{\v});}
  \foreach \v in {0.50,0.60,0.70,0.80,0.90}{
    \node[anchor=east, font=\ticklab, text=quietink] at (-0.14,\Y{\v}) {\v};}

  % --- the chance rule. Drawn, and named, because the axis is NIR. ----------
  \draw[ink, line width=0.5pt, densely dotted] (-0.10,\Y{0.50}) -- (4.32,\Y{0.50});
  \draw[ink, line width=0.5pt, densely dotted] ( 4.58,\Y{0.50}) -- (10.67,\Y{0.50});
  \draw[ink, line width=0.5pt, densely dotted] (10.93,\Y{0.50}) -- (\R,\Y{0.50});
  \node[anchor=east, font=\ticklab\itshape, text=ink] at (\R,\Y{0.483}) {chance};

  % --- what the axis is -----------------------------------------------------
  \node[anchor=west, font=\subhead, text=quietink]
    at (-0.64,\Y{0.933}) {$\NIR$, \Rframe{} frame};

  % --- column dividers and column heads -------------------------------------
  \draw[rulegrey, line width=0.4pt] (\Da,-0.05) -- (\Da,1.05);
  \draw[rulegrey, line width=0.4pt] (\Db,-0.05) -- (\Db,1.05);
  \node[anchor=south, align=center, text=ink] at (2.24,1.07)
    {{\headlab all scored conditions}\\[0.2ex]{\subhead\color{quietink}$n = \num{1394}$}};
  \node[anchor=south, align=center, text=ink] at (7.62,1.07)
    {{\headlab common support}\\[0.2ex]{\subhead\color{quietink}$n = \num{1192}$}};
  \node[anchor=south, align=center, text=ink] at (11.91,1.07)
    {{\headlab MoA support}\\[0.2ex]{\subhead\color{quietink}$n = \num{218}$}};

  % =========================================================================
  % COLUMN 1 -- all scored conditions, n = 1394
  % =========================================================================
  % Four fanned leaders. Dots are dodged sideways only where values coincide:
  % control_copy 0.506, random 0.501, generic 0.500 are 0.8mm and 0.16mm apart
  % vertically here, against a 1.0mm dot.
  \lead{0.30}{0.537}{1.20}{0.645}{accent}
  \lead{0.30}{0.506}{1.20}{0.583}{ink}
  \lead{0.58}{0.501}{1.20}{0.532}{ink}
  \lead{0.86}{0.500}{1.20}{0.478}{ink}

  \dotat{0.30}{0.854}{ink}
  \dotat{0.30}{0.537}{accent}
  \dotat{0.30}{0.506}{ink}
  \dotat{0.58}{0.501}{ink}
  \dotat{0.86}{0.500}{ink}

  \labat{0.52}{0.854}{ink}{within-well reference}
  \labat{1.20}{0.645}{accent}{model}
  \labat{1.20}{0.583}{ink}{\ttfamily control\_copy}
  \labat{1.20}{0.532}{ink}{\ttfamily random}
  \labat{1.20}{0.478}{ink}{\ttfamily generic}

  % =========================================================================
  % COLUMN 2 -- common support of the drug lookups, n = 1192
  % =========================================================================
  % Coverage of the reference, (arm - 0.50) / (reference - 0.50), both on THIS
  % support. Braces first so the dots' halos are not cut by a brace stem.
  \draw[decorate, decoration={brace, amplitude=3.5pt}, accent, line width=0.4pt]
    (5.95,\Y{0.500}) -- (5.95,\Y{0.549});
  \node[anchor=east, font=\marklab, text=accent] at (5.76,\Y{0.5245}) {$13.9\%$};
  \draw[decorate, decoration={brace, amplitude=3.5pt}, ink, line width=0.4pt]
    (6.60,\Y{0.500}) -- (6.60,\Y{0.913});
  \node[anchor=east, font=\marklab, text=ink] at (6.42,\Y{0.710}) {$115.7\%$};

  \dotat{7.15}{0.913}{ink}
  \dotat{7.15}{0.857}{ink}
  \dotat{7.15}{0.822}{ink}
  \dotat{7.15}{0.549}{accent}

  \labat{7.37}{0.913}{ink}{\ttfamily drug\_lookup}
  \labat{7.37}{0.857}{ink}{within-well reference}
  \labat{7.37}{0.822}{ink}{\ttfamily drug\_lookup\_1}
  \labat{7.37}{0.549}{accent}{model}

  % =========================================================================
  % COLUMN 3 -- conditions with an MoA partner pool, n = 218
  % =========================================================================
  \dotat{11.12}{0.596}{accent}
  \dotat{11.12}{0.552}{ink}

  \labat{11.34}{0.596}{accent}{model}
  \labat{11.34}{0.552}{ink}{\ttfamily moa\_lookup}

\end{tikzpicture}

\caption[The baseline ladder]{\textbf{A predictor that needs no model, no prompt
and no checkpoint scores above the reference the model is measured against, and
the three columns may not be read across.} Marks are $\NIR$ in the \Rframe{}
frame, where chance is $0.50$. Each column is scored on its own support
($n = \num{1394}$, $\num{1192}$ and $\num{218}$), so a mark in one may not be set
beside a mark in another; even the reference differs between the first two. On
the common support the model covers $13.9\%$ of the distance from chance to the
within-well split-half precision reference, and a training-only per-drug lookup
covers $115.7\%$ of it, which is a property of the reference and not headroom
(\cref{tab:lookup}). On the MoA support the model scores $0.596$ against
\texttt{moa\_lookup}'s $0.552$, a $+0.0436$ the interval does not separate from
zero. No interval is drawn: the source gives them on differences between arms,
not on the level of one arm, and those are in \cref{tab:lookup}. \emph{Reading
the labels}: monospace names an arm exactly as \cref{tab:lookup} names it, roman
names the reference in prose, the accent marks the model, and there is no fourth
style. A label a rule would strike or a neighbour would crowd is lifted clear on
a hairline; the dot carries the value. Horizontal position inside a column means
nothing, and arms whose scores nearly coincide are dodged sideways so each keeps
its own dot.}
\label{fig:ladder}
\end{figure}

% ---------------------------------------------------------------------------
\begin{table}[htbp]
\begin{fullwidth}
\begin{threeparttable}
\caption{\textbf{A training-only per-drug lookup is the bar to beat.} Reference,
model and drug-agnostic arms use all $\num{1394}$ conditions; the drug lookups
exist on a common support of $\num{1192}$, where model and reference score
$0.549$ and $0.857$. There, model minus \texttt{drug\_lookup} is
\ci{-0.3631}{-0.3982}{-0.3280}; the lookup scores $0.913$ pooled, $0.980$ on
trained conditions and $0.890$ on held-out wells, where its same-well advantage
is absent. \texttt{moa\_lookup} has narrower support ($n = \num{218}$) and is not
separated from the model (\ci{+0.0436}{-0.0480}{+0.1352}); its pooled row must
not be compared with rows on other supports.}
\label{tab:lookup}
\small
\begin{tabularx}{\linewidth}{@{}%
  >{\raggedright\arraybackslash}p{0.28\linewidth}
  >{\raggedright\arraybackslash}p{0.22\linewidth}
  S[table-format=1.3]
  >{\raggedright\arraybackslash}X@{}}
\toprule
\thead{Arm} & \thead{What it knows} & {\thead{$\NIR$}} & \thead{Note} \\
\midrule
within-well split-half reference & disjoint cell halves, one well & 0.854
  & precision, not a replicate; $0.857$ on the lookups' support \\
\addlinespace
\texttt{drug\_lookup} (training-only) & drug's mean residual, other lines & 0.913
  & the bar to beat; $n = \num{1192}$ \\
\texttt{drug\_lookup\_1} & the same, from one other cell line & 0.822
  & still beats the model; same support \\
\texttt{moa\_lookup} ($n = \num{218}$) & drugs sharing this drug's MoA & 0.552
  & narrower support; not separated \\
\addlinespace
model & control cell $+$ drug & 0.537
  & $0.549$ on the lookups' $n = \num{1192}$ \\
\addlinespace
\texttt{control\_copy} & nothing (copies the control) & 0.506 & drug-agnostic \\
\texttt{generic} & nothing (predicts the average) & 0.500 & the zero vector \\
\texttt{random} & nothing & 0.501 & floor \\
\bottomrule
\end{tabularx}
\begin{tablenotes}[flushleft]\footnotesize
\item Rows on different supports are not comparable. Chance is $0.50$ in the
\Rframe{} frame under the exchangeability condition audited in
\cref{sec:res-metrics}.
\end{tablenotes}
\end{threeparttable}
\end{fullwidth}
\end{table}

% ---------------------------------------------------------------------------
\paragraph{The pre-registered middle branch fired} Of the three branches written
before the result, the middle one fired. The model reads the drug and recovers
only a small fraction of its residual; transfer to held-out combinations is
\notestablished (\cref{sec:res-generalisation}). Whether its sensitivity to the
prompt reduces to the \texttt{Mechanism:} field is not settled here. On the
$\num{218}$ conditions where that comparison exists, model $-$
\texttt{moa\_lookup} $=
\ci{+0.0436}{-0.0480}{+0.1352}$, clustered on cell line and well, an interval
spanning zero.

\paragraph{This is DrEval's finding, in transcriptomic space} The finding travels
furthest of anything here, needing no checkpoint, no prompt and no cell-sentence
encoding. DrEval reports that deep learning models barely outperform a naive
predictor of mean drug and cell-line effects~\cite{dreval}, and our
\texttt{drug\_lookup} is
the $\beta(d)$ term of that additive model, measured where the
$\textrm{generic}(c)$ term has already been
subtracted.\seealso{\cref{sec:res-transfer} measures what that constant leaves
out} On a $946$-gene transcriptomic response, not a scalar potency, that drug
main effect reaches $0.913$, above the $0.857$ reference itself. DrEval's
inherited warning~\cite{partin}, that good pair-holdout performance is reachable
by predicting each drug's training average, is here the hypothesis. It holds.

\paragraph{The missing head-to-head, and why the bar survives it} A third
objection is that this chapter's model is the wrong one. A conditional model
built for perturbation response~\cite{cpa,scgen,gears} was never run on this
split, so the head-to-head a reader wants is missing, and that is a real gap. But
the ladder is architecture-free, so a competitor can be dropped in and scored,
and what it has to clear for a seen drug in a new context is
\texttt{drug\_lookup}, $0.913$ pooled, or $0.890$ on the held-out split. Not
$0.500$. For an unseen drug no lookup exists, so that regime needs a different
bar and gets one in \cref{sec:res-scope}.

\begin{scopelimit}[carried into \cref{sec:res-scope} and \cref{sec:limitations}]
The two regimes must not share a number. \texttt{drug\_lookup} is the bar for a
\emph{seen} drug in a new context. No lookup exists for a drug that never
appeared in training, so nothing here sets a bar for that regime. Neither may
rows on different supports be combined: $\num{1394}$, $\num{1192}$ and
$\num{218}$ select conditions by different rules, and a distance to the reference
is comparable within a support and not between supports.
\end{scopelimit}

\paragraph{Why the lookup beating the reference is not headroom} On the common
support the lookup scores \emph{above} the reference,
$\texttt{drug\_lookup} - \textrm{reference} = +0.056$, tempting to read as ``the
modelling target has vanished''. Three reasons, all favouring the lookup, block
that inference. The reference is scored against a half-sample truth and every
baseline against the full-sample truth; \texttt{drug\_lookup} averages
$\approx 38$ conditions against that single noisy half; and $\NIR$ near $0.91$
compresses most of the cosine range into a few thousandths. The coverage figures
are themselves distances to a split-half reference \cref{sec:methods-data} calls
optimistic; that optimism sits in the denominator and if anything makes both
coverage figures too small, while the half-sample truth of the first reason
pushes the other way. The net magnitude is unknown and the ordering carries the
claim.

\paragraph{Averaging is not why the lookup beats the model} Of those three
reasons only the averaging could also deflate the comparison with the model, and
\texttt{drug\_lookup\_1} tests it by drawing on one other cell line. It scores
$0.822$, and paired on the same support the model trails it by
\ci{-0.2723}{-0.3074}{-0.2372}, clustered two-way on cell line and well. This
drug's mean residual in one other cell line beats a $1$B-parameter model handed
this condition's own control cells. Neither gap is an average over a mixed
population (\cref{fig:paired-lookup}): condition by condition, the model scores
above \texttt{drug\_lookup} on $85$ of the $\num{1192}$ and above
\texttt{drug\_lookup\_1} on $214$, $7.1\%$ and $18.0\%$, counts over conditions
carrying no interval.

\paragraph{The bar is not equally clean on both sides of the split} The pooled
$0.913$ covers two populations. On trained conditions the lookup scores $0.980$.
On held-out wells, where it cannot borrow the target's own well, it scores
$0.890$ against the model's $0.533$ on the \needsaudit{893} held-out-well
conditions where both are scored, a count this thesis quotes two ways and has
under audit; the paired gap there is \ci{-0.357}{-0.403}{-0.311}. Nothing in the
argument turns on the exact size of that population. The ordering holds on both
sides, and the held-out side, where the well advantage is absent, is the side a
claim about transfer has to survive.

% ---------------------------------------------------------------------------
% \begin{fullwidth}: two square scatters side by side, each with its own
% difference histogram beneath, is a 174mm object -- squeezed to the 142mm body
% block the two panels lose the aspect that makes the identity line readable as
% a diagonal. Drawn at 174mm by thesis_v2/figs/paired-lookup.py, so NO width=
% key: it is placed at scale 1.0 and its type prints at the size it was set in.
% Every number it draws is listed with its source pointer in
% figs/fig-paired-lookup.json.
\begin{figure}[htbp]
\begin{fullwidth}
\centering
\includegraphics{fig-paired-lookup}
\caption[The model beats the lookup on 7.1\% of conditions]{\textbf{The headline
gap is a tilt of nearly the whole population and not a mean over a mixed one:
the model scores above \texttt{drug\_lookup} on $7.1\%$ of conditions, and above
the single-cell-line lookup on $18.0\%$.} Every mark is one of the $\num{1192}$
conditions of the common support --- those on which a training-only lookup
exists --- and both axes are $\NIR$ in the \Rframe{} frame, comparison set other
drugs in the same cell line. On this support the model scores $0.549$ and the
within-well split-half precision reference $0.857$; the $\num{1394}$-condition
figures quoted elsewhere ($0.537$ and $0.854$) are a different population and
are not combined with these. \textbf{(a)}~Against \texttt{drug\_lookup}, the
drug's mean residual in all other cell lines, $0.913$ pooled: the model is above
the identity line on $85$ of $\num{1192}$ conditions, $7.1\%$.
\textbf{(b)}~Against \texttt{drug\_lookup\_1}, the same quantity from a single
other cell line, $0.822$ pooled: $214$ of $\num{1192}$, $18.0\%$. Marker shape
gives the split, $\num{299}$ \texttt{train} and \needsaudit{893}
\texttt{unseen\_combo}; no \texttt{unseen\_drug} condition appears, because no
training lookup exists for a drug held out of training, and that regime is
therefore not addressed by this figure at all. The lower panels are the
per-condition paired difference on a shared count axis, with the paired mean
drawn below the counts as a point and bar, the bar a $95\%$ bootstrap interval
clustered two-way on cell line and well: \ci{-0.3631}{-0.3982}{-0.3280} in (a)
and \ci{-0.2723}{-0.3074}{-0.2372} in (b). The win counts are counts over
conditions and carry no interval; ties ($23$ and $28$ conditions) count as not a
win. The $0.50$ rules are chance on both axes, labelled once in (a).}
\label{fig:paired-lookup}
\end{fullwidth}
\end{figure}

\begin{claim}
The honest summary is the paired gap on common support ($n = \num{1192}$),
\ci{-0.3631}{-0.3982}{-0.3280}. The model covers $13.9\%$ of the distance from
chance to the within-well split-half precision reference where a training-only
lookup covers all of it and more. Those are common-support figures, model $0.549$
against reference $0.857$; over all $\num{1394}$ scored conditions, not all of
which admit a lookup, the model is $0.537$ against a reference of $0.854$, and
the two populations are not combined. The pooled gap does combine the two sides
of the split, on which the lookup is not equally clean. On held-out conditions
alone it is \ci{-0.357}{-0.403}{-0.311}, and the lookup's $0.913$ is $0.980$ on
trained conditions against $0.890$ on held-out ones. That the lookup exceeds the
reference is a property of the reference, not a headroom estimate; the reference
is a single noisy half-sample while the lookup averages many cell lines, and
$\NIR$ compresses near its top. Nor is averaging what beats the model. Restricted
to a \emph{single} other cell line the lookup scores $0.822$ and the model still
trails it by \ci{-0.2723}{-0.3074}{-0.2372}. This is DrEval's finding replicated
in transcriptomic space, the drug main effect at $0.913$ where a $1$B-parameter
model reaches $0.549$. Being a bar and not a result is what makes it reusable,
and any method proposed for seen-drug transfer on this atlas has to clear it, not
merely chance.
\end{claim}

% ===========================================================================
\beatsection{R}{8}{Same-target retrieval is the strongest observed lead}{sec:res-scope}
\claimlede{Same-target retrieval is the strongest observed lead on its own
restricted support; because the three channel arms differ in support, drug
roster, partner count and annotation coverage, that is a descriptive result and
not an identified biological ranking.}

\begin{question}[8]
The bar just set exists only for a drug that appeared in training. A drug that
never did can be reached at all only through some property it shares with drugs
that were seen, a protein target, a mechanism, a chemical shape. The instrument
predicts each drug from its annotated neighbours in the same cell line, against a
null matched on everything but the channel itself. One channel carries a lead, on
a very short roster of drugs.
\end{question}

\noindent
For a channel $X$, a drug's residual is predicted as the mean residual of the
other drugs in the same cell line sharing property $X$ with it, the construction
of \texttt{moa\_lookup} (\cref{sec:res-lookup}) generalised from mechanism to any
annotated property.

\paragraph{The null has to differ in exactly one respect} The gate's estimand
requires the channel arm and its null to differ in the channel and in nothing
else. Every channel is paired with a plate-matched random null carrying the same
number of partner drugs, drawn from the same cell line, with the same split
across the target's own plate and other plates. Mechanism- and target-related
drugs are co-plated more often than chance, so a null matched on count alone
would differ in plate membership too, two differences where the estimand allows
one. Where a plate-matched null spans too few cell lines the channel is reported
\textsc{untestable}. ``We could not build the control'' must never be written as
``closed''. Coverage is printed first as well, because a channel annotated on few
drugs was never really tested and cannot be closed by a null result.

\paragraph{The partner pool was wrong, and every verdict here is from the re-run}
The gate was intended to draw partners only from drugs that were in training,
which is what makes it a fair contest in every split. An audit of its own
implementation, not of the already-restricted lookup arms of
\cref{sec:res-lookup}, found the pool spanning every retained condition in the
cell line. The restriction is now in the code. Of the
\needsaudit[what retained counts]{4{,}131} retained conditions only
\needsaudit{$\num{2523}$} are training, so nearly two in five the old pool could
draw on were ineligible.
Every verdict here is from the corrected run, which moved the target gap down,
the mechanism gap up, and chemistry's down by half.\corrmark{the channel gate's
partner pool, restricted and re-run}

% ---------------------------------------------------------------------------
\begin{table}[htbp]
\begin{fullwidth}
\begin{threeparttable}
\caption{\textbf{Coverage first, then the gap.} The drugs behind each channel are
stated before any verdict, because a null result on a channel never really tested
is not a closure. The plate-matched null is the estimand-correct control; the
count-matched column shows sensitivity to plate membership, not a causal share.
The gate asks the plate-matched lower bound to clear a pre-set margin of $0.03$
clustered on cell line and well. On that axis, and only there, protein target
clears it ($+0.0745$) while mechanism ($+0.0279$) and chemistry ($+0.0035$) do
not. All three exclude zero, so the two that fail are inconclusive rather than
equivalent to their nulls, and neither may be written as closed.}
\label{tab:channels}
\small
\begin{tabularx}{\linewidth}{@{}%
  >{\raggedright\arraybackslash}p{0.15\linewidth}
  >{\centering\arraybackslash}p{0.13\linewidth}
  >{\centering\arraybackslash}p{0.06\linewidth}
  >{\centering\arraybackslash}p{0.23\linewidth}
  >{\centering\arraybackslash}p{0.11\linewidth}
  >{\raggedright\arraybackslash}X@{}}
\toprule
\thead{Channel} & \thead{Drugs (annot./scored)} & \thead{$n$}
  & \thead{vs plate-matched null} & \thead{(vs count-matched)}
  & \thead{Verdict} \\
\midrule
\textbf{protein target} & 264 / 18 & 712 & \ci{+0.1351}{+0.0745}{+0.1958}
  & $+0.1498$ & live, cell-line axis$^{\dagger}$ \\
\addlinespace
mechanism (MoA) & 180 / 26 & 921 & \ci{+0.0805}{+0.0279}{+0.1332}
  & $+0.0623$ & inconclusive \\
\addlinespace
chemical structure & 377 / 94 & \needsaudit{4131}
  & \ci{+0.0261}{+0.0035}{+0.0487} & $+0.0122$ & inconclusive \\
\bottomrule
\end{tabularx}
\begin{tablenotes}[flushleft]\footnotesize
\item $^{\dagger}$The axis is written into the target verdict because ``live'' is
not usable without it. Clustered on the drug, the target interval falls short of
this margin against both estimand-correct nulls and clears it only against the
weaker count-matched one.
\item The counts are annotations in Tahoe's drug metadata and not testable
coverage here; within this cache only $18$, $26$ and $94$ of the $94$ scored
drugs received a channel arm at all.
\end{tablenotes}
\end{threeparttable}
\end{fullwidth}
\end{table}

% ---------------------------------------------------------------------------
% \input, not \includegraphics: figs/gate.tex is TikZ and carries its own float,
% fullwidth block, caption and \label{fig:gate}. Its twelve intervals are
% literals transcribed from RESULTS_cluster/ and listed again, to machine
% precision and with their n, cluster keys, df and source paths, in the sibling
% figs/fig-gate.json.
% ===========================================================================
% gate.tex -- the channel gate, twelve intervals in three blocks.
% \input by Sections/04d-baselines.tex, replacing the inline tikzpicture that
% sat at lines 391-455 there. Compiles against thesis_v2/preamble.tex and
% nothing else. NO external data at draw time: every coordinate below is a
% literal transcribed from RESULTS_cluster/ and listed again, to full machine
% precision and with its n and cluster keys, in the sibling figs/fig-gate.json.
%
% SOURCES, one per block. Nothing here is computed in the figure.
%   block 1  RESULTS_cluster/channel_gate_v4.json
%            channels.<ch>.{plate_matched,different_plate}.{gap,ci}
%            ci is the two-way cell-line x well interval (df 38 / 40 / 41).
%   block 2  RESULTS_cluster/channel_gate_v4_byaxis.json
%            channels.target.<null>.ci_drug_well (df 17); the gap is the same
%            number as in block 1 -- re-clustering moves the interval only.
%   block 3  RESULTS_cluster/channel_gate_heldout_targets.json
%            channels.<ch>.heldout_targets.{gap,ci}, cell line x well again
%            (df 4 / 10 / 39).
% Each row's n is the .n at that same path, transcribed with it; block 2 re-uses
% block 1's records, so its counts are block 1's and only the clustering differs.
% Every drawn value equals the four-decimal figure the prose and tab:channels
% already print; the agreement was checked value by value, in both directions,
% against 04d-baselines.tex lines 350-354, 369-376, 470-479, 521-526 and 684.
%
% WHAT CHANGED FROM THE VERSION IT REPLACES, and why. The old drawing was sound
% -- greyscale-safe, no chartjunk, blocks separated, caption forbidding a read
% across blocks -- and five things were wrong with it, the fifth found in a
% review of this file's own first version:
%
%   1. THE AXIS HAD NO NAME. It carried ticks, a dotted line and a dashed line
%      and nowhere said what was being measured, so a reader could not tell an
%      NIR gap from a DRF or a coverage without going back to tab:channels one
%      page earlier. There is now an axis title over the data area naming the
%      quantity (a gap in NIR), its direction (arm minus matched null), its
%      scale (NIR runs 0 to 1, chance 0.50) and what zero means. This is the
%      one thing the figure genuinely needed. It ran to two lines as first
%      written; the third, naming the frame, is fix 5 below.
%
%   2. THE ORIGIN WAS LABELLED IN WORDS while the other ticks were numerals --
%      "zero" beside 0.10 and 0.20 on one axis. The origin is now 0.00, in the
%      same face, size and colour as its neighbours, on the dotted rule.
%
%   3. THE TWO CALLOUTS SAT ON TOP OF EACH OTHER. "zero" and "margin" labelled
%      two rules 5.5mm apart, in two different styles, at the same height. The
%      zero rule now carries a numeral in the tick row; the margin keeps a
%      worded callout, raised so the gap between the two label boxes is
%      0.82 y-units = 4.4mm of clear space (y=5.4mm; inner sep is pinned at
%      1pt below so this is box arithmetic, not judgement, and the ink gap is
%      wider still). The dashed rule is carried up to meet its own callout, so
%      no separate leader is needed.
%
%   4. THE AXIS STOPPED AT ZERO but two rows did not. mechanism-on-four runs to
%      -0.0846 and chemistry-different-plate to -0.0097, so rows extended into
%      a region with no tick and no label. The tick set now runs
%      -0.10, 0.00, 0.10, 0.20 at a uniform 0.10 pitch and the negative region
%      the third block occupies is labelled like any other.
%
%   5. THE AXIS NAMED THE QUANTITY BUT NOT THE FRAME, and no row carried its
%      own n. Fix 1 said what was measured, a gap in NIR, and left out where it
%      was measured. This document runs two measurement frames that both put
%      chance at 0.50, and conflating them is its most common misread -- the
%      running head exists for that reason alone -- so an NIR figure naming
%      neither the frame nor the set the score is ranked against is
%      under-labelled in exactly the way fix 1 was, one level up. A third title
%      line now names both, in the wording 04d-baselines.tex uses at 714-716
%      and fig:comparator uses in its caption: the residual frame, scored
%      against the other drugs in the same cell line. The caption says it too,
%      so the figure does not depend on the running head to carry its frame.
%      Twelve rows also carried no counts while tab:channels, one page earlier,
%      carries three; each row now prints its own n in a column of its own,
%      headed n, so what a block rests on is visible where the block is. The
%      three plate-matched counts in block 1 are tab:channels' n column, value
%      for value; the chemistry count keeps that table's audit flag, since a
%      quantity flagged in one place and printed clean in another is exactly
%      the divergence the flag exists to prevent. Every count is in
%      fig-gate.json, per row, with its source path.
%
% GEOMETRY, all of it derived rather than judged. x=108mm is the DATA area,
% spanning \alo..\ahi = -0.115..0.265, so 0.10 of NIR is 28.4mm.
%   padding    set by the data: leftmost ink is the -0.0846 whisker cap, 8.7mm
%              clear of the left edge; rightmost the +0.2460 cap, 5.7mm clear
%              of the right.
%   labels     column 60.0mm = 0.5556 x-units, set by the longest string in it,
%              the block head "targets restricted to drugs held out of
%              training" at 166.38pt = 58.5mm (measured, figs/
%              _measure_gate_strings.tex). Nothing in the column reaches the
%              data area, so no label can foul the -0.10 rule, which stands
%              4.3mm inside the left edge. Longest ROW label is 42.2mm.
%   counts     the n column is right-aligned inside that same 60.0mm, its right
%              edge 1.5mm clear of the data area and therefore 5.8mm clear of
%              the -0.10 rule. Its widest entry is the audit-flagged 4,131 chip
%              at 24.60pt = 8.65mm, and that entry sits on the row carrying the
%              longest label, "chemical structure, plate-matched" at 120.04pt =
%              42.19mm; label and count are 6.96mm apart with the 1pt inner sep
%              counted at both ends, which is the tightest pair in the column.
%              The column takes space the label column already had, so the
%              drawing does not widen: the harness reports the same 484.15pt
%              before and after.
%   rules      block separators start 1mm left of the label column and end at
%              the right edge of the data area.
%   width      1.5648 x-units = 169.0mm against the 174mm a fullwidth float may
%              use; 5mm of slack. MEASURED by figs/_harness_gate.tex, which
%              \typeout s \wd of the tikzpicture against the real preamble.
%   title      all three lines fit inside the 108mm data area they are centred
%              on (78.76mm, 88.08mm and 102.13mm measured), so none widens the
%              drawing. The block sits 0.65 y-units = 3.51mm higher than it did
%              to make room for the frame line; height is the only dimension
%              this change moves, and the float is nowhere near \topfraction.
%
% TYPE. \scriptsize throughout, matching fig:clustering and fig:confound. This
% drawing is not scaled by anything, so 8pt here prints as 8pt.
% ===========================================================================
\begin{figure}[htbp]
\begin{fullwidth}
\centering
\begin{tikzpicture}[x=108mm, y=5.4mm, font=\small, inner sep=1pt]
  \def\alo{-0.115}\def\ahi{0.265}  % data range of the x axis
  \def\lx{-0.5556}                 % left edge of the row-label column (60.0mm)
  \def\rx{-0.5648}                 % left end of the block rules      (61.0mm)
  \def\ytop{0.72}                  % top of the tick rules and the zero rule
  \def\ybot{-16.1}                 % bottom of the same
  \def\cx{-1.5mm}                  % right edge of the n column, from x=0
  \newcommand{\vx}[1]{{(#1-\alo)/(\ahi-\alo)}}

  % --- one interval row: #1 y, #2 lo, #3 estimate, #4 hi, #5 label, #6 n ---
  % FIX 5: #6 is the row's own count, right-aligned against \cx. A column, not a
  % caption list: twelve counts in prose cannot be matched to twelve rows by eye,
  % and this is how tab:channels carries them one page earlier.
  \newcommand{\gaprow}[6]{%
    \draw[quietink, line width=0.9pt] (\vx{#2},#1) -- (\vx{#4},#1);
    \draw[quietink, line width=0.5pt]
      (\vx{#2},#1-0.22) -- (\vx{#2},#1+0.22);
    \draw[quietink, line width=0.5pt]
      (\vx{#4},#1-0.22) -- (\vx{#4},#1+0.22);
    \fill[ink] (\vx{#3},#1) circle (1.5pt);
    \node[anchor=west, font=\scriptsize, text=ink] at (\lx,#1) {#5};
    \node[anchor=east, font=\scriptsize, text=quietink, xshift=\cx]
      at (0,#1) {#6};}
  \newcommand{\blockhead}[2]{%
    \node[anchor=west, font=\scriptsize\bfseries, text=accent]
      at (\lx,#1) {#2};}

  % --- FIX 1: the axis title. What is measured, in which direction, on what
  % scale, and what zero means. Centred on the data area, not on the drawing.
  % FIX 5 adds the middle line: in which frame, and against which set the score
  % is ranked. Both frames in this document put chance at 0.50, so the scale
  % line below does not distinguish them and cannot be left to do that work.
  \node[anchor=south, font=\scriptsize\bfseries, text=ink] at (0.5,4.70)
    {Gap in normalised inverse rank, channel arm $-$ matched null};
  \node[anchor=south, font=\scriptsize, text=quietink] at (0.5,4.05)
    {residual frame, $\NIR$ scored against the other drugs in the same
     cell line};
  \node[anchor=south, font=\scriptsize, text=quietink] at (0.5,3.40)
    {$\NIR$ runs $0$ to $1$ with chance at $0.50$; a gap of $0.00$
     is no advantage from the channel};

  % --- the scale. FIX 4: the tick set now covers the negative region. ------
  \foreach \v/\t in {-0.10/$-0.10$, 0.10/0.10, 0.20/0.20}{
    \draw[rulegrey!55, line width=0.3pt] (\vx{\v},\ytop) -- (\vx{\v},\ybot);
    \node[anchor=south, font=\scriptsize, text=quietink]
      at (\vx{\v},\ytop+0.06) {\t};}

  % zero, and FIX 2: a numeral, in the same row and register as the others.
  \draw[ink, line width=0.5pt, densely dotted] (\vx{0},\ytop) -- (\vx{0},\ybot);
  \node[anchor=south, font=\scriptsize, text=quietink]
    at (\vx{0},\ytop+0.06) {0.00};

  % the pre-set relevance margin. FIX 3: the rule is carried up to its own
  % callout, 4.4mm of clear space above the tick numerals.
  \draw[accent, line width=0.6pt, dashed] (\vx{0.03},2.35) -- (\vx{0.03},\ybot);
  % (the 1.3mm stand-off is an xshift, not an addition to \vx: \vx expands to
  % a braced pgfmath group and "{...}+0.012" is not a legal coordinate.)
  \node[anchor=west, font=\scriptsize, text=accent, xshift=1.3mm]
    at (\vx{0.03},2.35) {pre-set relevance margin, $0.03$};

  % --- block 1: cell line x well, the axis the gate was run on -------------
  % The n column is headed once, on the first block head's own row: that head is
  % 38.4mm of the 60.0mm column and stops 11mm short of the column, so the two
  % cannot touch. Blocks 2 and 3 keep the same column position and the caption
  % names it, so the head is not repeated over rules it would be read across.
  \blockhead{0}{clustered on cell line and well}
  \node[anchor=east, font=\scriptsize, text=quietink, xshift=\cx] at (0,0) {$n$};
  \gaprow{-1}{+0.0745}{+0.1351}{+0.1958}{protein target, plate-matched}{\num{712}}
  \gaprow{-2}{+0.0648}{+0.1322}{+0.1996}{protein target, different-plate}{\num{599}}
  \gaprow{-3}{+0.0279}{+0.0805}{+0.1332}{mechanism, plate-matched}{\num{921}}
  \gaprow{-4}{+0.0243}{+0.0746}{+0.1249}{mechanism, different-plate}{\num{781}}
  % the flagged count is tab:channels' own \needsaudit{4131}, carried here so the
  % same quantity is not hedged in the table and clean in the figure.
  \gaprow{-5}{+0.0035}{+0.0261}{+0.0487}{chemical structure, plate-matched}%
    {\needsaudit{4{,}131}}
  \gaprow{-6}{-0.0097}{+0.0118}{+0.0333}{chemical structure, different-plate}%
    {\num{4090}}

  % --- block 2: the same target arm, re-clustered on the drug --------------
  % Same records as rows -1, -2 and the count-matched arm, so the counts repeat:
  % re-clustering changes df (38 to 17), never n.
  \draw[rulegrey, line width=0.3pt] (\rx,-7.0) -- (1.00,-7.0);
  \blockhead{-7.8}{the same arm, clustered on drug and well}
  \gaprow{-8.8}{+0.0287}{+0.1351}{+0.2416}{protein target, plate-matched}{\num{712}}
  \gaprow{-9.8}{+0.0199}{+0.1322}{+0.2445}{protein target, different-plate}{\num{599}}
  \gaprow{-10.8}{+0.0536}{+0.1498}{+0.2460}{protein target, count-matched}{\num{712}}

  % --- block 3: targets restricted to drugs held out of training -----------
  % Here the row label counts DRUGS and the column counts conditions; the two
  % are different numbers for the same row, which is why the column earns its
  % place most in this block.
  \draw[rulegrey, line width=0.3pt] (\rx,-11.8) -- (1.00,-11.8);
  \blockhead{-12.6}{targets restricted to drugs held out of training}
  \gaprow{-13.6}{+0.0499}{+0.1197}{+0.1894}{protein target, on two drugs}{\num{161}}
  \gaprow{-14.6}{-0.0846}{+0.0439}{+0.1723}{mechanism, on four}{\num{216}}
  \gaprow{-15.6}{-0.0029}{+0.0378}{+0.0785}{chemical structure, on fourteen}{\num{898}}
\end{tikzpicture}
\caption[The channel gate, by conditioning]{\textbf{The gate's verdict is decided
by what the interval is conditioned on, not by where the point estimate falls.}
Every row is residual-frame $\NIR$ scored against the other drugs in the same
cell line: the arm predicts a drug's residual from its annotated neighbours
there, and its matched null draws the same number of partner drugs from that
cell line.
The dot is the estimate, the bar its $95\%$ two-way cluster-robust interval, the
dashed line the pre-set relevance margin of $0.03$ that the gate asks the
\emph{lower} bound to clear. The count beside each row label is the number of
conditions that row scores, not of drugs; the first block's plate-matched counts
are the $n$ of \cref{tab:channels}. The second block re-clusters the same target
arm on the drug, the unit a claim about unseen drugs has to generalise across,
leaving the point estimates untouched; the lower bound now falls short against both
estimand-correct nulls and only the weaker count-matched one clears it. The third
restricts the targets to drugs held out of training. Rows may be read within a
block and not across blocks.}
\label{fig:gate}
\end{fullwidth}
\end{figure}


% ---------------------------------------------------------------------------
\paragraph{The population the question is actually about} The repair restricted
where the \emph{predictors} come from, not which conditions are \emph{scored}.
The targets are all retained conditions, spanning $94$ drugs, of which $14$ are
held out of training entirely, the $10$ of the \texttt{unseen\_drug} quota plus
four appearing only in held-out wells. The target arm covers $18$ drugs,
\textbf{two} of them held out. Restrict the targets to held-out drugs and,
clustered two-way on cell line and well as everything else here is, the target
gap is \ci{+0.1197}{+0.0499}{+0.1894} on those two; mechanism falls from
$+0.0805$ to \ci{+0.0439}{-0.0846}{+0.1723} on four and chemistry to
\ci{+0.0378}{-0.0029}{+0.0785} on fourteen, both spanning zero
(\cref{fig:gate}, third block). Only the target channel survives, and it
survives on two drugs.

\paragraph{What eighteen drugs are made of} The target arm's $18$ annotated drugs
are almonertinib, amsacrine, cepharanthine, clonidine, elagolix, gemfibrozil,
goserelin, idarubicin, lumateperone, mitoxantrone, neratinib, norepinephrine,
olanzapine, rapamycin, relugolix, temsirolimus, tofacitinib and tofacitinib
citrate. Their mean partner count is $1.19$. Of the arm's $712$ predictions,
$579$ come from a \emph{single} partner drug and $133$ from two, against $2.46$
partners for mechanism and $5.00$ for chemical similarity. And sometimes that
partner is an unusually close relative. \texttt{Tofacitinib} and
\texttt{Tofacitinib (citrate)}, the free base and the citrate salt of one drug,
are scored in the same $20$ cell lines and have one partner each; so are
rapamycin and its ester prodrug temsirolimus, in $22$ lines. Neratinib and
almonertinib pair the same way in $20$ lines, and those two are chemically
distinct, so the pattern is not only duplicate chemistry. The channel
establishes that a same-target neighbour predicts the drug. That a target
\emph{annotation} generalises across a class is the stronger statement; testing
it would mean dropping the nearest analogue from the pool, and that was not run.

\paragraph{The two drugs the verdict rests on} They are \textbf{amsacrine} and
\textbf{relugolix}, and because the duplicate-chemistry pattern above would
otherwise be the obvious objection to them, their pools are worth naming.
Amsacrine is annotated to \texttt{TOP2A} and its partners are doxorubicin,
epirubicin, idarubicin and mitoxantrone; relugolix is annotated to
\texttt{GNRHR} and its partners are elagolix and goserelin. Neither pool
contains a salt, prodrug or stereoisomer of the drug it predicts, so this
verdict is not the tofacitinib case in another guise --- though amsacrine's
pool is three anthracyclines and a close structural analogue of them, which is
a drug class and not a duplicate. The support is $\num{161}$ conditions over
$39$ cell lines and $5$ treated wells, $70$ of them amsacrine in two wells and
$91$ relugolix in three; the five wells are the smaller cluster count and are
where this interval's $t_4$ reference comes from.\seealso{\cref{sec:methods-data}:
why the smaller cluster count sets the reference distribution}

\begin{scopelimit}[carried into \cref{sec:res-field} and \cref{sec:limitations}]
What this section measures is \emph{channel availability} --- whether the
information a channel carries is present in the atlas at all --- and not
unseen-drug predictive performance, which two drugs cannot establish for the
target channel and which mechanism does not survive at all. Every verdict here,
including those printed in \cref{tab:channels}, is to be read in those terms. The
bound is on similarity transfer through each channel and not on any model, which
could combine channels or learn a non-smooth structure-to-response map that a
nearest-neighbour average cannot express, though at a few hundred drugs the
lookup is effectively the bound.
\end{scopelimit}

\paragraph{Why the clustering axis decides this section} Two conditions from one
treated well share a physical assignment and a batch; two from one cell line
share the biology. Every interval on a mean in this chapter is therefore two-way
cluster-robust,
\begin{equation}
 V \;=\; V_{g_1} \;+\; V_{g_2} \;-\; V_{\textrm{intersection}},
 \label{eq:cgm}
\end{equation}
where $g_1$ and $g_2$ are the cell line and the treated well, whose intersection
is the single condition, so $V_{\textrm{intersection}}$ is the ordinary
independent-sampling variance (\cref{sec:methods-data}). Such a variance is
asymptotic in the \emph{number of clusters} and not in the number of conditions,
so the reference distribution is Student's $t$ on $\min(G_{g_1},G_{g_2}) - 1$.
Wells sit inside drugs rather than crossing them, so on the drug axis the
intersection is the well and not the condition; retaining the crossed form there
subtracts less and widens the interval, and the drug-axis figures in this chapter
are the demanding ones by construction.

\paragraph{And on this section's own arm it changes the verdict} The gate
clusters on cell line and treated well, right for a statement about this
experiment. But the claim the channel is asked to support, that an \emph{unseen
drug} is reachable through its target annotation, generalises over drugs, and
only $18$ annotated drugs carry the target arm. Against the plate-matched null
the target gap is \ci{+0.1351}{+0.0287}{+0.2416} clustered on drug and well
($\mathrm{df} = 17$; \cref{fig:gate}, second block against the first), whose
lower bound falls $0.0013$ short of the margin;
against the different-plate null it is \ci{+0.1322}{+0.0199}{+0.2445}. Both
exclude zero; neither clears $0.03$. Only the count-matched null clears it here,
\ci{+0.1498}{+0.0536}{+0.2460}, and that is the control already set aside as
weaker. Mechanism, which excludes zero on the cell-line axis, spans it on the
drug axis in all three constructions. The honest summary is therefore narrower
than \emph{live}, and the axis is chosen by the claim, not by which interval is
narrower.

\paragraph{Sensitivity to plate matching, measured and not argued away} Mechanism
partners are co-plated with the target $0.404$ of the time against $0.362$ for a
random draw; elsewhere the diagnostic is negligible, protein target $0.368$
against $0.366$ and chemical similarity $0.370$ against $0.367$. Under a third
construction drawing all partners from other plates, target and mechanism keep
intervals clear of zero and chemistry's spans it
(\cref{tab:fieldattribution}). Absolute $\NIR$ there is $0.638$, $0.575$ and
$0.499$ against nulls of $0.506$, $0.500$ and $0.487$, so removing the shared
plate does not push the arms below chance. It barely moves target or mechanism,
two and seven per cent off their plate-matched gaps, while chemistry attenuates
again. But it also changes support and available partners, so the attenuation
does not identify co-plating as its cause.

\paragraph{The ordering is descriptive, and one of its terms is a floor} Target,
then mechanism, then chemistry, under every construction and again on held-out
targets. The target and chemistry intervals do not overlap under any of the three
nulls; each adjacent pair's do, and on the drug axis even the two ends overlap.
These arms differ in support, drug set, partner count and annotation coverage, so
the order does not identify a biological ranking. One correction runs in a known
direction. The residual frame subtracts a plate-local generic, and
\cref{sec:methods-residual} shows that for a co-plated mechanism class part of
that class's own shared programme goes with it, so mechanism's middle position is
a floor and not an estimate. Below target the two failures differ. Mechanism's
lower bound falls $0.0021$ short, a knife-edge that is not a negative result in
either direction; chemistry's interval is the \emph{narrowest} of the three,
$0.045$ wide against $0.121$ for target, but straddles the margin. Neither is a
closure.

\paragraph{A pre-registration this overturns} Two design documents concluded that
unseen-drug generalisation was out of reach, and both inferred it from the
chemical-structure channel alone, Morgan fingerprints and a molecular language
model. An earlier drug-knowledge specification had instead identified protein
targets as ``the key feature'', since transcriptional response is driven by which
protein a drug hits, something structure encodes only indirectly. The
same-target lead shows that generalisation from chemistry to every channel was
not licensed. The gate agrees with the specification about which lead to pursue;
with arms on different supports it does not validate it as a hierarchy. Tahoe's
own metadata carries protein-target annotations for $264$ drugs ($280$ distinct
symbols, $550$ edges), which no analysis in this project had read.

\paragraph{Three caveats bound the target verdict} The first is annotation
coverage. The arm covers $18$ of the $94$ scored drugs. The second is scale. A
shared property buys a fraction of what the drug itself buys; the target arm sits
$+0.14$ above chance where a same-drug lookup for a seen drug sits $+0.41$
(\cref{sec:res-lookup}, a different condition set, so a comparison of scale and
not a paired one). The third is overlap. Whether target adds anything over MoA is
unsettled here; settling it would need a leave-one-mechanism-out split, since a
held-out drug's mechanism class is generally still represented in training.

\begin{claim}
The target arm excludes zero on every construction under both clusterings, and on
held-out targets, the population the unseen-drug question is actually about, the
gap is \ci{+0.1197}{+0.0499}{+0.1894} on two drugs, clustered on cell line and
well as everything else here is. On the full arm it clears the pre-set relevance
margin of $0.03$ on the cell-line axis under all three constructions; clustered
on the drug, the axis a statement about unseen drugs rides on, it clears the
margin only against the weaker count-matched null and falls short against both
estimand-correct ones. The arm rests on $18$ annotated drugs, two of them held
out. Mechanism excludes zero on the cell-line axis, falls $0.0021$ short of the
margin there, spans zero on the drug axis, and spans zero again on held-out
targets (\ci{+0.0439}{-0.0846}{+0.1723}); it cannot carry an unseen-drug claim in
any conditioning, and the frame understates rather than flatters it. Chemistry's
interval is the narrowest of the three but straddles the margin, so it does not
locate the effect wholly either side of that threshold; it does rule out a large
effect, and under the different-plate construction it no longer clears zero,
though support and partner availability change with plate separation, so that
attenuation does not identify co-plating as its cause. None of the three may be
written as closed, and none as live without naming the axis. Similarity transfer
through the target channel is measurable and is the lead to pursue first; whether
it is large enough to matter is unresolved at $18$ drugs, and mechanism is
unresolved and not excluded, graded in a frame that subtracts part of its own
signal before scoring it.
\end{claim}

% ===========================================================================
\beatsection{ER}{9}{Which part of the prompt is read remains unresolved}{sec:res-field}

\begin{question}[9]
The previous section asked whether the information an unseen drug would need is
present in the atlas at all. The complementary question is whether the model
reads what it is already handed, and which part of it, the name or the mechanism
line beside it. The instrument swaps one span of the prompt at a time and leaves
every other byte identical. It does not settle the question, and what follows is
why, and what would settle it.
\end{question}

\noindent
The chapter has established that the prediction responds to that information as a
block; it has not established \emph{which} part of it is read. Saying so
precisely matters, because the next step writes itself. Append the protein
targets to the prompt, retrain, and see whether unseen-drug performance rises.
That arm assumes the bottleneck is information availability. Everything in this
chapter points the other way, but the direct evidence is weaker than earlier
drafts claimed.

\paragraph{The instrument, and the one guard it needs} The scramble has until now
replaced the drug name and the mechanism together, so a non-null gap establishes
sensitivity to the combined prompt without attributing it to either span. A field
decomposition swaps exactly one span to the \texttt{opposite} partner, the name
with the mechanism kept, the mechanism with the name kept, and both together. One
guard is essential. If the swap partner happens to share the original's
mechanism, swapping only the mechanism produces a prompt identical to the model's
own and the arm scores exactly zero by construction, which read naively is
indistinguishable from ``the model ignores the mechanism'', the very claim under
test. Such conditions are dropped, which is why the mechanism arm carries fewer
conditions.

\paragraph{What the arms said, and why it does not stand} On the earlier target
build the decomposed arms suggested the prompt sensitivity is carried almost
entirely by the name span. Swapping only the drug name gave $\gap$ of
\ci{+0.0809}{+0.0592}{+0.1010} ($n = 570$); swapping only the mechanism gave
$+0.0091$ [$-0.0174$, $+0.0383$] ($n = 325$); both together gave $+0.1024$
[$+0.0827$, $+0.1236$]. All were scored against the \texttt{opposite}
comparator that \cref{sec:res-generalisation} shows inflates every gap read
against it. They \emph{were} re-run under the corrected evaluation
(\cref{tab:fieldattribution}, Panel B), but the defect survives it, being a
property of the swap partner and not of the target build. The claim that the
model ``reads the name and not the mechanism'' is therefore withdrawn pending
comparators whose own scores are empirically near chance. What the current run
establishes is sensitivity to the combined prompt, with the mechanism field's
contribution unresolved.\corrmark{the name-over-mechanism reading, withdrawn
pending neutral comparators}

% ---------------------------------------------------------------------------
\begin{table}[htbp]
\begin{fullwidth}
\begin{threeparttable}
\caption{\textbf{Availability does not establish use.} Panel A shows target and
mechanism retrieval surviving the different-plate null while chemistry does not;
only target clears the relevance margin, and only on the cell-line axis, on an
arm covering $18$ drugs (\cref{sec:res-scope}). Intervals are two-way
cluster-robust over cell line and treated well. Panel B cannot attribute the
combined-prompt effect between name and mechanism, because its two partners lie
on opposite sides of chance ($0.434$ and $0.520$), inflating one gap and
suppressing the other. Comparator-free margins reverse the apparent ordering
($+0.0367$ for name, $+0.0506$ for mechanism), so the split is unresolved pending
comparators whose own scores are empirically near chance. Identical mechanism
swaps are omitted, which is why that row has fewer conditions.}
\label{tab:fieldattribution}
\small
\begin{tabularx}{\linewidth}{@{}%
  >{\raggedright\arraybackslash}p{0.20\linewidth}
  >{\centering\arraybackslash}p{0.15\linewidth}
  >{\centering\arraybackslash}p{0.08\linewidth}
  >{\raggedright\arraybackslash}X@{}}
\toprule
\multicolumn{4}{@{}l}{\thead{Panel A: channel availability --- metadata-neighbour
retrieval}} \\
\thead{Channel} & \thead{Drugs (annot./scored)} & \thead{$n$}
  & \thead{Gap vs plate-matched random null} \\
\midrule
protein target & $264$ / $18$ & $712$ & \ci{+0.1351}{+0.0745}{+0.1958} \\
mechanism (MoA) & $180$ / $26$ & $921$ & \ci{+0.0805}{+0.0279}{+0.1332} \\
chemical structure & $377$ / $94$ & \needsaudit{4131} & \ci{+0.0261}{+0.0035}{+0.0487} \\
\midrule
\multicolumn{4}{@{}l}{the same three, against the stricter
\textsc{different-plate} null} \\
protein target & & $599$ & \ci{+0.1322}{+0.0648}{+0.1996} \\
mechanism (MoA) & & $781$ & \ci{+0.0746}{+0.0243}{+0.1249} \\
\textbf{chemical structure} & & $4090$ & \ci{+0.0118}{-0.0097}{+0.0333} \\
\bottomrule
\end{tabularx}

\vspace{0.8em}

\begin{tabularx}{\linewidth}{@{}%
  >{\raggedright\arraybackslash}p{0.25\linewidth}
  >{\centering\arraybackslash}p{0.27\linewidth}
  >{\centering\arraybackslash}p{0.08\linewidth}
  >{\raggedright\arraybackslash}X@{}}
\toprule
\multicolumn{4}{@{}l}{\thead{Panel B: field use --- one span swapped, every other
byte identical}} \\
\thead{Span swapped} & \thead{$\gap$} & \thead{$n$}
  & \thead{Comparator direction (partner $\NIR$)} \\
\midrule
drug name, mechanism kept & \ci{+0.1026}{+0.0625}{+0.1427} & $1394$
  & inflated ($0.434$) \\
mechanism, drug name kept & \ci{+0.0303}{-0.0078}{+0.0685} & $738$
  & suppressed ($0.520$) \\
both spans (the standard arm) & \ci{+0.1139}{+0.0781}{+0.1496} & $1394$
  & inflated ($0.423$) \\
\bottomrule
\end{tabularx}
\begin{tablenotes}[flushleft]\footnotesize
\item Panel A is repeated from \cref{tab:channels} so the two panels can be read
against each other; it is the same measurement, not a second one.
\end{tablenotes}
\end{threeparttable}
\end{fullwidth}
\end{table}

% ---------------------------------------------------------------------------
\paragraph{Why the two panels cannot be added together} Panel A says a channel
carries information in the atlas; Panel B asks whether the model uses what it is
handed. Neither answers the other's question, and an ordering that reverses when
the comparator is removed is not an attribution.

\paragraph{The practical consequence, narrower than earlier drafts drew} The
target-prompt arm is not retired by measurement. It is untested. What argues for
building it carefully is the shape of the failure above. The model's deficit is
not confined to missing information about unseen drugs, because on seen drugs it
covers $13.9\%$ of the distance from chance to the within-well reference, against
a training-only lookup that covers all of it and more (\cref{sec:res-lookup}).
What argues for building it at all is the one thing a lookup cannot do, handle a
drug it has never seen, the regime where \cref{sec:res-scope} finds measurable
transfer through the protein-target channel, on $18$ drugs.

\paragraph{The opportunity is not the part that fails to transfer}
\cref{sec:res-transfer} puts $1 - \Tcoef$ at \ci{44.7\%}{40.7\%}{48.6\%}, and it
is tempting to call that the model's headroom. It is not. That figure is a
shortfall in cross-context similarity, mixing cell line, dose, well and
estimator; it is neither a variance share nor an interaction. With $286$ of $287$
treatments in exactly one well there is no replicate from which to separate a
context-specific remainder from measurement noise, which is why the coefficient
is a disattenuated cross-line cosine and not a fitted variance
decomposition.\seealso{\cref{sec:res-transfer} for the estimator and its five
checks} Whether any of that shortfall is learnable structure and not noise is
separated nowhere here. The opportunity is the unseen drug, and whether the model
can be made to take it is not answered in this chapter.

% ===========================================================================
\beatsection{ER}{0}{The investigation in one table}{sec:res-summary}

\begin{question}[0]
The results are distributed across two measurement frames, and both put chance in
the same place. That coincidence is the trap this chapter has to leave its reader
safe from. What follows is a ledger of the principal comparison arms, frame by
frame, then the results that are estimates and belong outside it.
\end{question}

\noindent
The expression frame of \cref{tab:allarms} scores a predicted cell profile
against other drugs on the same plate. The residual frame scores a predicted
residual against other drugs in the same cell line, after the generic programme
has been subtracted.

% ---------------------------------------------------------------------------
% MOVED HERE, ahead of tab:allarms -- PLACEMENT ONLY, nothing in the figure or
% its caption touched. It was declared after the ledger table, and the two
% floats are introduced from the same page: the table took the only top slot on
% the page after it and this figure was pushed one further, landing two pages
% from the \Cref that introduces it. Declared first it takes that slot instead.
% The table loses nothing by taking the later one -- it is \cref'd on BOTH sides
% of where it now sets (the frame-defining paragraph above, and "What the ledger
% deliberately leaves out" below), so it still lands one page from a reference,
% while this figure has exactly one \Cref and no second anchor to fall back on.
% MEASURED: figure 2 pages from its \Cref -> 1; table still one page from its
% nearest reference and still heading a text page, not a page of its own; the
% three lines that close the chapter are still not stranded (the regression the
% note below was written against); page count unchanged by this move.
\begin{figure}[htbp]
% NOT \begin{fullwidth}. Laid out at the 142mm body block: the argument is a
% differential stretch between two rulers, and widening the block widens both
% equally, so the extra 32mm buys nothing the reader can use. The tikzpicture
% lives in figs/two-scales.tex and every number it draws, with its source
% pointer, is listed in figs/two-scales.json.
\centering
% ===========================================================================
% two-scales.tex -- v2. fig:two-scales. GENERATED by two-scales_build.py from
% RESULTS_cluster/nir_sameplate.json and RESULTS_cluster/re_v3.json. Do not edit
% by hand: re-run the builder. Every number is listed in two-scales.json.
%
% BODY WIDTH, 142 mm exactly. Not \begin{fullwidth}. Contains a tikzpicture and
% nothing else, so a section \input's it inside its own figure float -- the
% pattern of figs/comparators.tex and figs/frames.tex, not of fig-frame.tex.
% The float wrapper and caption are at the foot of this file, commented out,
% ready to paste into Sections/04d-baselines.tex.
%
% THE CONSTRUCTION. Chance (0.50) is drawn at the same x in both scales and each
% frame's OWN within-well split-half precision reference is drawn at the same x
% as the other's. Nothing else is aligned. Because the two chance-to-reference
% distances are 0.0764 and 0.3539 of NIR, the two rulers come out 4.63 times
% apart, and that differential stretch is the whole argument. The calibration
% bars, flush to the right margin, make it measurable rather than merely felt.
%
% REGISTER, inherited from fig:ladder: mono for the identifiers that name a
% column in the results JSON (\texttt{generic}, \texttt{random},
% \texttt{control\_copy}, \texttt{orth}, \texttt{drug\_lookup}); roman for arms
% the thesis names in prose (model, linear map, control-copy, global mean,
% within-well reference). Accent is reserved for the model and for chance.
%
% WHAT IS DELIBERATELY ABSENT: every interval (no arm here is quoted with one at
% the pooled level, so all markers are bare); the identifiable-stratum
% expression rows 0.768 / 0.766, which are a different support; and the three
% supports of fig:ladder. The one derived quantity drawn, +0.039, appears ONLY
% as the thing being struck out.
% ===========================================================================
\begin{tikzpicture}[
  x=1cm, y=1cm,
  scaleline/.style={ink, line width=0.7pt, line cap=round},
  guide/.style={rulegrey, densely dotted, line width=0.3pt},
  chanceguide/.style={accent!65, densely dotted, line width=0.45pt},
  leader/.style={rulegrey!85, line width=0.25pt},
  headrule/.style={rulegrey!60, line width=0.3pt},
  hardrule/.style={ink, line width=1.0pt},
  connector/.style={ink!60, densely dashed, line width=0.7pt},
  calib/.style={quietink, line width=0.9pt, line cap=butt},
  fh/.style={font=\sffamily\footnotesize\bfseries, text=accent, inner sep=0pt},
  fhs/.style={font=\footnotesize, text=quietink, inner sep=0pt},
  gh/.style={anchor=north, font=\sffamily\scriptsize, text=accent,
             align=center, inner sep=1pt},
  val/.style={font=\footnotesize, text=ink, align=center, inner sep=1.5pt},
  arm/.style={font=\scriptsize, text=quietink, inner sep=1.5pt},
  armM/.style={font=\footnotesize\bfseries, text=accent, inner sep=1.5pt},
  note/.style={font=\scriptsize, text=quietink, align=flush left, inner sep=0pt},
  noteR/.style={font=\scriptsize, text=quietink, align=flush right, inner sep=0pt},
]

% ---- the bounding box IS the body block; nothing may stick out of it
\path (0.0000,1.60) rectangle (14.2000,11.3000);

% ================ the two alignment anchors
\node[gh] at (3.6000,11.3000) {chance $=0.50$\\ the one value the two frames share};
\node[gh] at (10.6000,11.3000) {within-well split-half reference\\ each frame's own, drawn at the same place};
\draw[chanceguide] (3.6000,10.5000) -- (3.6000,3.2500);
\draw[guide]       (10.6000,10.5000) -- (10.6000,3.2500);

% ================ UPPER SCALE -- expression frame
\node[fh,  anchor=west] at (0.0000,10.30) {expression frame};
\node[fhs, anchor=east] at (14.2000,10.30) {within plate, tier 2, $n=\num{606}$};
\draw[headrule] (0.0300,10.00) -- (14.1700,10.00);
\draw[scaleline] (2.35,8.8500) -- (12.30,8.8500);
\draw[scaleline] (1.15,8.8500) -- (1.80,8.8500);   % the broken stub
\draw[ink, line width=0.5pt] (1.96,8.7200) -- (2.10,8.9800);
\draw[ink, line width=0.5pt] (2.08,8.7200) -- (2.22,8.9800);
\fill[quietink] (1.4750,8.8500) circle (1.4pt);   % global mean, off scale
\fill[quietink] (3.5605,8.8500) circle (1.4pt);   % linear map
\fill[quietink] (3.9220,8.8500) circle (1.4pt);   % control-copy
\fill[ink]      (10.6000,8.8500) circle (1.6pt);   % within-well reference
\fill[accent]   (3.3904,8.8500) circle (1.9pt);   % model
\draw[accent!85, line width=0.5pt] (3.6000,8.7400) -- (3.6000,8.9600);
\node[val, anchor=south] at (1.4750,9.0000) {global mean\\ $0.180$};
\node[val, anchor=south] at (10.6000,9.0000) {$0.576$};
\draw[leader] (3.3904,8.7800) -- (3.20,8.5300);
\node[armM, anchor=east] at (3.16,8.5300) {model\;$0.498$};
\draw[leader] (3.5605,8.7800) -- (3.20,8.1800);
\node[arm, anchor=east] at (3.16,8.1800) {linear map\;$0.500$};
\draw[leader] (3.9220,8.7800) -- (3.20,7.8300);
\node[arm, anchor=east] at (3.16,7.8300) {control-copy\;$0.504$};
\node[note, anchor=north west, text width=3.30cm] at (0.00,7.58) {Scale broken at the left: $0.180$ lies $4.2\times$ this frame's chance-to-reference distance below chance.};
\draw[calib] (9.5900,7.10) -- (14.1700,7.10);
\draw[calib] (9.5900,7.02) -- (9.5900,7.18);
\draw[calib] (14.1700,7.02) -- (14.1700,7.18);
\node[noteR, anchor=north east] at (14.2000,6.96) {$0.05$ of $\NIR$ on this scale};

% ================ the subtraction the figure exists to forbid
\draw[white, line width=2.2pt] (3.3904,8.7400) -- (3.5864,7.7500);
\draw[connector] (3.3904,8.7400) -- (3.8555,6.2400);
\draw[connector] (3.9232,5.8600) -- (4.3259,3.7100);
\node[draw=ink, line width=0.5pt, rounded corners=1.5pt, fill=panelbg, inner sep=0.17cm, text width=4.70cm, anchor=south west, font=\scriptsize, text=quietink, align=flush left] (nosub) at (4.3000,6.2200) {\tikz[baseline=(eq.base)]{ \node[inner sep=0pt, font=\footnotesize, text=ink, anchor=west] (eq) at (0.50,0) {$0.537 - 0.498 = +0.039$}; \draw[ink, line width=1.6pt] ([xshift=-0.10cm]eq.west) -- ([xshift=0.10cm]eq.east); \fill[ink] (0.15,0.075) circle (0.130); \draw[white, line width=0.9pt, line cap=round] (0.085,0.010) -- (0.215,0.140); \draw[white, line width=0.9pt, line cap=round] (0.085,0.140) -- (0.215,0.010); }\\[3pt] is the subtraction the shared chance value invites, and it is not a quantity: two different frames, measured against two different references.};
\node[fill=panelbg, inner xsep=3pt, inner ysep=1pt, anchor=west, font=\sffamily\scriptsize\bfseries, text=ink] at ([xshift=0.30cm]nosub.north west) {Do not subtract};

% ================ the scope limit, drawn as a rule and not as a caption
\draw[hardrule] (0.0300,6.0500) -- (14.1700,6.0500);
\node[anchor=north west, font=\sffamily\scriptsize, text=ink, align=flush left, inner sep=0pt, text width=9.90cm] at (4.30,5.9200) {\textbf{Scope limit.} No number in one frame may be set beside a number in the other.};
% the connector is stopped at the rule, with the panel's own mark
\fill[ink] (3.8894,6.0500) circle (0.190);
\draw[white, line width=1.0pt, line cap=round] (3.7944,5.9550) -- (3.9844,6.1450);
\draw[white, line width=1.0pt, line cap=round] (3.7944,6.1450) -- (3.9844,5.9550);

% ================ LOWER SCALE -- residual frame
\node[fh,  anchor=west] at (0.0000,5.05) {residual frame};
\node[fhs, anchor=east] at (14.2000,5.05) {cell-line sets, $n=\num{1394}$};
\draw[headrule] (0.0300,4.75) -- (14.1700,4.75);
\draw[scaleline] (2.35,3.6000) -- (12.30,3.6000);
\fill[quietink] (3.6000,3.6000) circle (1.4pt);   % generic
\fill[quietink] (3.6113,3.6000) circle (1.4pt);   % scramble_orth
\fill[quietink] (3.6126,3.6000) circle (1.4pt);   % random
\fill[quietink] (3.7118,3.6000) circle (1.4pt);   % control_copy
\fill[ink]    (10.6000,3.6000) circle (1.6pt);   % within-well reference
\fill[ink]    (11.7597,3.6000) circle (1.6pt);   % drug_lookup
\fill[accent] (4.3259,3.6000) circle (1.9pt);   % model
\draw[accent!85, line width=0.5pt] (3.6000,3.4900) -- (3.6000,3.7100);
\draw[leader] (4.3259,3.6900) -- (4.46,3.8400);
\node[val, anchor=south west, align=left, text=accent, font=\footnotesize\bfseries] at (4.44,3.7800) {model\\ $0.537$};
\node[val, anchor=south] at (10.6000,3.7500) {$0.854$};
\draw[leader] (11.7597,3.6900) -- (11.98,3.8400);
\node[arm, anchor=south west, align=left, text=ink] at (11.96,3.7800) {\texttt{drug\_lookup}\\ $0.913$};
\draw[leader] (3.6818,3.5100) -- (3.44,3.26);
\draw[leader] (3.40,3.26) -- (3.40,2.29);
\draw[leader] (3.40,3.2600) -- (3.34,3.2600);
\node[arm, anchor=east] at (3.30,3.2600) {\texttt{generic}\;$0.500$};
\draw[leader] (3.40,2.9400) -- (3.34,2.9400);
\node[arm, anchor=east] at (3.30,2.9400) {scramble, \texttt{orth}\;$0.501$};
\draw[leader] (3.40,2.6200) -- (3.34,2.6200);
\node[arm, anchor=east] at (3.30,2.6200) {\texttt{random}\;$0.501$};
\draw[leader] (3.40,2.3000) -- (3.34,2.3000);
\node[arm, anchor=east] at (3.30,2.3000) {\texttt{control\_copy}\;$0.506$};
\draw[leader] (4.3659,3.5100) -- (4.66,3.30);
\node[note, anchor=north west, text width=5.50cm] at (4.70,3.36) {The $13.9\%$ margin this chapter quotes for the model (\cref{fig:ladder}) is computed on the $n=\num{1192}$ common support, where the model is $0.549$ against a reference of $0.857$. It is not a measurement of the $n=\num{1394}$ gap drawn here.};
\draw[calib] (13.1810,3.00) -- (14.1700,3.00);
\draw[calib] (13.1810,2.92) -- (13.1810,3.08);
\draw[calib] (14.1700,2.92) -- (14.1700,3.08);
\node[noteR, anchor=north east] at (14.2000,2.86) {$0.05$ of $\NIR$ on this scale};
\node[noteR, anchor=north east, text width=3.85cm] at (14.2000,2.44) {The two rulers differ by $4.6\times$.};
\end{tikzpicture}
%
% ---------------------------------------------------------------------------
% READY TO PASTE into Sections/04d-baselines.tex, at sec:res-summary, just after
% the paragraph "The same chance value, two different references".
%
% \begin{figure}[htbp]
% \centering
% % ===========================================================================
% two-scales.tex -- v2. fig:two-scales. GENERATED by two-scales_build.py from
% RESULTS_cluster/nir_sameplate.json and RESULTS_cluster/re_v3.json. Do not edit
% by hand: re-run the builder. Every number is listed in two-scales.json.
%
% BODY WIDTH, 142 mm exactly. Not \begin{fullwidth}. Contains a tikzpicture and
% nothing else, so a section \input's it inside its own figure float -- the
% pattern of figs/comparators.tex and figs/frames.tex, not of fig-frame.tex.
% The float wrapper and caption are at the foot of this file, commented out,
% ready to paste into Sections/04d-baselines.tex.
%
% THE CONSTRUCTION. Chance (0.50) is drawn at the same x in both scales and each
% frame's OWN within-well split-half precision reference is drawn at the same x
% as the other's. Nothing else is aligned. Because the two chance-to-reference
% distances are 0.0764 and 0.3539 of NIR, the two rulers come out 4.63 times
% apart, and that differential stretch is the whole argument. The calibration
% bars, flush to the right margin, make it measurable rather than merely felt.
%
% REGISTER, inherited from fig:ladder: mono for the identifiers that name a
% column in the results JSON (\texttt{generic}, \texttt{random},
% \texttt{control\_copy}, \texttt{orth}, \texttt{drug\_lookup}); roman for arms
% the thesis names in prose (model, linear map, control-copy, global mean,
% within-well reference). Accent is reserved for the model and for chance.
%
% WHAT IS DELIBERATELY ABSENT: every interval (no arm here is quoted with one at
% the pooled level, so all markers are bare); the identifiable-stratum
% expression rows 0.768 / 0.766, which are a different support; and the three
% supports of fig:ladder. The one derived quantity drawn, +0.039, appears ONLY
% as the thing being struck out.
% ===========================================================================
\begin{tikzpicture}[
  x=1cm, y=1cm,
  scaleline/.style={ink, line width=0.7pt, line cap=round},
  guide/.style={rulegrey, densely dotted, line width=0.3pt},
  chanceguide/.style={accent!65, densely dotted, line width=0.45pt},
  leader/.style={rulegrey!85, line width=0.25pt},
  headrule/.style={rulegrey!60, line width=0.3pt},
  hardrule/.style={ink, line width=1.0pt},
  connector/.style={ink!60, densely dashed, line width=0.7pt},
  calib/.style={quietink, line width=0.9pt, line cap=butt},
  fh/.style={font=\sffamily\footnotesize\bfseries, text=accent, inner sep=0pt},
  fhs/.style={font=\footnotesize, text=quietink, inner sep=0pt},
  gh/.style={anchor=north, font=\sffamily\scriptsize, text=accent,
             align=center, inner sep=1pt},
  val/.style={font=\footnotesize, text=ink, align=center, inner sep=1.5pt},
  arm/.style={font=\scriptsize, text=quietink, inner sep=1.5pt},
  armM/.style={font=\footnotesize\bfseries, text=accent, inner sep=1.5pt},
  note/.style={font=\scriptsize, text=quietink, align=flush left, inner sep=0pt},
  noteR/.style={font=\scriptsize, text=quietink, align=flush right, inner sep=0pt},
]

% ---- the bounding box IS the body block; nothing may stick out of it
\path (0.0000,1.60) rectangle (14.2000,11.3000);

% ================ the two alignment anchors
\node[gh] at (3.6000,11.3000) {chance $=0.50$\\ the one value the two frames share};
\node[gh] at (10.6000,11.3000) {within-well split-half reference\\ each frame's own, drawn at the same place};
\draw[chanceguide] (3.6000,10.5000) -- (3.6000,3.2500);
\draw[guide]       (10.6000,10.5000) -- (10.6000,3.2500);

% ================ UPPER SCALE -- expression frame
\node[fh,  anchor=west] at (0.0000,10.30) {expression frame};
\node[fhs, anchor=east] at (14.2000,10.30) {within plate, tier 2, $n=\num{606}$};
\draw[headrule] (0.0300,10.00) -- (14.1700,10.00);
\draw[scaleline] (2.35,8.8500) -- (12.30,8.8500);
\draw[scaleline] (1.15,8.8500) -- (1.80,8.8500);   % the broken stub
\draw[ink, line width=0.5pt] (1.96,8.7200) -- (2.10,8.9800);
\draw[ink, line width=0.5pt] (2.08,8.7200) -- (2.22,8.9800);
\fill[quietink] (1.4750,8.8500) circle (1.4pt);   % global mean, off scale
\fill[quietink] (3.5605,8.8500) circle (1.4pt);   % linear map
\fill[quietink] (3.9220,8.8500) circle (1.4pt);   % control-copy
\fill[ink]      (10.6000,8.8500) circle (1.6pt);   % within-well reference
\fill[accent]   (3.3904,8.8500) circle (1.9pt);   % model
\draw[accent!85, line width=0.5pt] (3.6000,8.7400) -- (3.6000,8.9600);
\node[val, anchor=south] at (1.4750,9.0000) {global mean\\ $0.180$};
\node[val, anchor=south] at (10.6000,9.0000) {$0.576$};
\draw[leader] (3.3904,8.7800) -- (3.20,8.5300);
\node[armM, anchor=east] at (3.16,8.5300) {model\;$0.498$};
\draw[leader] (3.5605,8.7800) -- (3.20,8.1800);
\node[arm, anchor=east] at (3.16,8.1800) {linear map\;$0.500$};
\draw[leader] (3.9220,8.7800) -- (3.20,7.8300);
\node[arm, anchor=east] at (3.16,7.8300) {control-copy\;$0.504$};
\node[note, anchor=north west, text width=3.30cm] at (0.00,7.58) {Scale broken at the left: $0.180$ lies $4.2\times$ this frame's chance-to-reference distance below chance.};
\draw[calib] (9.5900,7.10) -- (14.1700,7.10);
\draw[calib] (9.5900,7.02) -- (9.5900,7.18);
\draw[calib] (14.1700,7.02) -- (14.1700,7.18);
\node[noteR, anchor=north east] at (14.2000,6.96) {$0.05$ of $\NIR$ on this scale};

% ================ the subtraction the figure exists to forbid
\draw[white, line width=2.2pt] (3.3904,8.7400) -- (3.5864,7.7500);
\draw[connector] (3.3904,8.7400) -- (3.8555,6.2400);
\draw[connector] (3.9232,5.8600) -- (4.3259,3.7100);
\node[draw=ink, line width=0.5pt, rounded corners=1.5pt, fill=panelbg, inner sep=0.17cm, text width=4.70cm, anchor=south west, font=\scriptsize, text=quietink, align=flush left] (nosub) at (4.3000,6.2200) {\tikz[baseline=(eq.base)]{ \node[inner sep=0pt, font=\footnotesize, text=ink, anchor=west] (eq) at (0.50,0) {$0.537 - 0.498 = +0.039$}; \draw[ink, line width=1.6pt] ([xshift=-0.10cm]eq.west) -- ([xshift=0.10cm]eq.east); \fill[ink] (0.15,0.075) circle (0.130); \draw[white, line width=0.9pt, line cap=round] (0.085,0.010) -- (0.215,0.140); \draw[white, line width=0.9pt, line cap=round] (0.085,0.140) -- (0.215,0.010); }\\[3pt] is the subtraction the shared chance value invites, and it is not a quantity: two different frames, measured against two different references.};
\node[fill=panelbg, inner xsep=3pt, inner ysep=1pt, anchor=west, font=\sffamily\scriptsize\bfseries, text=ink] at ([xshift=0.30cm]nosub.north west) {Do not subtract};

% ================ the scope limit, drawn as a rule and not as a caption
\draw[hardrule] (0.0300,6.0500) -- (14.1700,6.0500);
\node[anchor=north west, font=\sffamily\scriptsize, text=ink, align=flush left, inner sep=0pt, text width=9.90cm] at (4.30,5.9200) {\textbf{Scope limit.} No number in one frame may be set beside a number in the other.};
% the connector is stopped at the rule, with the panel's own mark
\fill[ink] (3.8894,6.0500) circle (0.190);
\draw[white, line width=1.0pt, line cap=round] (3.7944,5.9550) -- (3.9844,6.1450);
\draw[white, line width=1.0pt, line cap=round] (3.7944,6.1450) -- (3.9844,5.9550);

% ================ LOWER SCALE -- residual frame
\node[fh,  anchor=west] at (0.0000,5.05) {residual frame};
\node[fhs, anchor=east] at (14.2000,5.05) {cell-line sets, $n=\num{1394}$};
\draw[headrule] (0.0300,4.75) -- (14.1700,4.75);
\draw[scaleline] (2.35,3.6000) -- (12.30,3.6000);
\fill[quietink] (3.6000,3.6000) circle (1.4pt);   % generic
\fill[quietink] (3.6113,3.6000) circle (1.4pt);   % scramble_orth
\fill[quietink] (3.6126,3.6000) circle (1.4pt);   % random
\fill[quietink] (3.7118,3.6000) circle (1.4pt);   % control_copy
\fill[ink]    (10.6000,3.6000) circle (1.6pt);   % within-well reference
\fill[ink]    (11.7597,3.6000) circle (1.6pt);   % drug_lookup
\fill[accent] (4.3259,3.6000) circle (1.9pt);   % model
\draw[accent!85, line width=0.5pt] (3.6000,3.4900) -- (3.6000,3.7100);
\draw[leader] (4.3259,3.6900) -- (4.46,3.8400);
\node[val, anchor=south west, align=left, text=accent, font=\footnotesize\bfseries] at (4.44,3.7800) {model\\ $0.537$};
\node[val, anchor=south] at (10.6000,3.7500) {$0.854$};
\draw[leader] (11.7597,3.6900) -- (11.98,3.8400);
\node[arm, anchor=south west, align=left, text=ink] at (11.96,3.7800) {\texttt{drug\_lookup}\\ $0.913$};
\draw[leader] (3.6818,3.5100) -- (3.44,3.26);
\draw[leader] (3.40,3.26) -- (3.40,2.29);
\draw[leader] (3.40,3.2600) -- (3.34,3.2600);
\node[arm, anchor=east] at (3.30,3.2600) {\texttt{generic}\;$0.500$};
\draw[leader] (3.40,2.9400) -- (3.34,2.9400);
\node[arm, anchor=east] at (3.30,2.9400) {scramble, \texttt{orth}\;$0.501$};
\draw[leader] (3.40,2.6200) -- (3.34,2.6200);
\node[arm, anchor=east] at (3.30,2.6200) {\texttt{random}\;$0.501$};
\draw[leader] (3.40,2.3000) -- (3.34,2.3000);
\node[arm, anchor=east] at (3.30,2.3000) {\texttt{control\_copy}\;$0.506$};
\draw[leader] (4.3659,3.5100) -- (4.66,3.30);
\node[note, anchor=north west, text width=5.50cm] at (4.70,3.36) {The $13.9\%$ margin this chapter quotes for the model (\cref{fig:ladder}) is computed on the $n=\num{1192}$ common support, where the model is $0.549$ against a reference of $0.857$. It is not a measurement of the $n=\num{1394}$ gap drawn here.};
\draw[calib] (13.1810,3.00) -- (14.1700,3.00);
\draw[calib] (13.1810,2.92) -- (13.1810,3.08);
\draw[calib] (14.1700,2.92) -- (14.1700,3.08);
\node[noteR, anchor=north east] at (14.2000,2.86) {$0.05$ of $\NIR$ on this scale};
\node[noteR, anchor=north east, text width=3.85cm] at (14.2000,2.44) {The two rulers differ by $4.6\times$.};
\end{tikzpicture}
%
% ---------------------------------------------------------------------------
% READY TO PASTE into Sections/04d-baselines.tex, at sec:res-summary, just after
% the paragraph "The same chance value, two different references".
%
% \begin{figure}[htbp]
% \centering
% % ===========================================================================
% two-scales.tex -- v2. fig:two-scales. GENERATED by two-scales_build.py from
% RESULTS_cluster/nir_sameplate.json and RESULTS_cluster/re_v3.json. Do not edit
% by hand: re-run the builder. Every number is listed in two-scales.json.
%
% BODY WIDTH, 142 mm exactly. Not \begin{fullwidth}. Contains a tikzpicture and
% nothing else, so a section \input's it inside its own figure float -- the
% pattern of figs/comparators.tex and figs/frames.tex, not of fig-frame.tex.
% The float wrapper and caption are at the foot of this file, commented out,
% ready to paste into Sections/04d-baselines.tex.
%
% THE CONSTRUCTION. Chance (0.50) is drawn at the same x in both scales and each
% frame's OWN within-well split-half precision reference is drawn at the same x
% as the other's. Nothing else is aligned. Because the two chance-to-reference
% distances are 0.0764 and 0.3539 of NIR, the two rulers come out 4.63 times
% apart, and that differential stretch is the whole argument. The calibration
% bars, flush to the right margin, make it measurable rather than merely felt.
%
% REGISTER, inherited from fig:ladder: mono for the identifiers that name a
% column in the results JSON (\texttt{generic}, \texttt{random},
% \texttt{control\_copy}, \texttt{orth}, \texttt{drug\_lookup}); roman for arms
% the thesis names in prose (model, linear map, control-copy, global mean,
% within-well reference). Accent is reserved for the model and for chance.
%
% WHAT IS DELIBERATELY ABSENT: every interval (no arm here is quoted with one at
% the pooled level, so all markers are bare); the identifiable-stratum
% expression rows 0.768 / 0.766, which are a different support; and the three
% supports of fig:ladder. The one derived quantity drawn, +0.039, appears ONLY
% as the thing being struck out.
% ===========================================================================
\begin{tikzpicture}[
  x=1cm, y=1cm,
  scaleline/.style={ink, line width=0.7pt, line cap=round},
  guide/.style={rulegrey, densely dotted, line width=0.3pt},
  chanceguide/.style={accent!65, densely dotted, line width=0.45pt},
  leader/.style={rulegrey!85, line width=0.25pt},
  headrule/.style={rulegrey!60, line width=0.3pt},
  hardrule/.style={ink, line width=1.0pt},
  connector/.style={ink!60, densely dashed, line width=0.7pt},
  calib/.style={quietink, line width=0.9pt, line cap=butt},
  fh/.style={font=\sffamily\footnotesize\bfseries, text=accent, inner sep=0pt},
  fhs/.style={font=\footnotesize, text=quietink, inner sep=0pt},
  gh/.style={anchor=north, font=\sffamily\scriptsize, text=accent,
             align=center, inner sep=1pt},
  val/.style={font=\footnotesize, text=ink, align=center, inner sep=1.5pt},
  arm/.style={font=\scriptsize, text=quietink, inner sep=1.5pt},
  armM/.style={font=\footnotesize\bfseries, text=accent, inner sep=1.5pt},
  note/.style={font=\scriptsize, text=quietink, align=flush left, inner sep=0pt},
  noteR/.style={font=\scriptsize, text=quietink, align=flush right, inner sep=0pt},
]

% ---- the bounding box IS the body block; nothing may stick out of it
\path (0.0000,1.60) rectangle (14.2000,11.3000);

% ================ the two alignment anchors
\node[gh] at (3.6000,11.3000) {chance $=0.50$\\ the one value the two frames share};
\node[gh] at (10.6000,11.3000) {within-well split-half reference\\ each frame's own, drawn at the same place};
\draw[chanceguide] (3.6000,10.5000) -- (3.6000,3.2500);
\draw[guide]       (10.6000,10.5000) -- (10.6000,3.2500);

% ================ UPPER SCALE -- expression frame
\node[fh,  anchor=west] at (0.0000,10.30) {expression frame};
\node[fhs, anchor=east] at (14.2000,10.30) {within plate, tier 2, $n=\num{606}$};
\draw[headrule] (0.0300,10.00) -- (14.1700,10.00);
\draw[scaleline] (2.35,8.8500) -- (12.30,8.8500);
\draw[scaleline] (1.15,8.8500) -- (1.80,8.8500);   % the broken stub
\draw[ink, line width=0.5pt] (1.96,8.7200) -- (2.10,8.9800);
\draw[ink, line width=0.5pt] (2.08,8.7200) -- (2.22,8.9800);
\fill[quietink] (1.4750,8.8500) circle (1.4pt);   % global mean, off scale
\fill[quietink] (3.5605,8.8500) circle (1.4pt);   % linear map
\fill[quietink] (3.9220,8.8500) circle (1.4pt);   % control-copy
\fill[ink]      (10.6000,8.8500) circle (1.6pt);   % within-well reference
\fill[accent]   (3.3904,8.8500) circle (1.9pt);   % model
\draw[accent!85, line width=0.5pt] (3.6000,8.7400) -- (3.6000,8.9600);
\node[val, anchor=south] at (1.4750,9.0000) {global mean\\ $0.180$};
\node[val, anchor=south] at (10.6000,9.0000) {$0.576$};
\draw[leader] (3.3904,8.7800) -- (3.20,8.5300);
\node[armM, anchor=east] at (3.16,8.5300) {model\;$0.498$};
\draw[leader] (3.5605,8.7800) -- (3.20,8.1800);
\node[arm, anchor=east] at (3.16,8.1800) {linear map\;$0.500$};
\draw[leader] (3.9220,8.7800) -- (3.20,7.8300);
\node[arm, anchor=east] at (3.16,7.8300) {control-copy\;$0.504$};
\node[note, anchor=north west, text width=3.30cm] at (0.00,7.58) {Scale broken at the left: $0.180$ lies $4.2\times$ this frame's chance-to-reference distance below chance.};
\draw[calib] (9.5900,7.10) -- (14.1700,7.10);
\draw[calib] (9.5900,7.02) -- (9.5900,7.18);
\draw[calib] (14.1700,7.02) -- (14.1700,7.18);
\node[noteR, anchor=north east] at (14.2000,6.96) {$0.05$ of $\NIR$ on this scale};

% ================ the subtraction the figure exists to forbid
\draw[white, line width=2.2pt] (3.3904,8.7400) -- (3.5864,7.7500);
\draw[connector] (3.3904,8.7400) -- (3.8555,6.2400);
\draw[connector] (3.9232,5.8600) -- (4.3259,3.7100);
\node[draw=ink, line width=0.5pt, rounded corners=1.5pt, fill=panelbg, inner sep=0.17cm, text width=4.70cm, anchor=south west, font=\scriptsize, text=quietink, align=flush left] (nosub) at (4.3000,6.2200) {\tikz[baseline=(eq.base)]{ \node[inner sep=0pt, font=\footnotesize, text=ink, anchor=west] (eq) at (0.50,0) {$0.537 - 0.498 = +0.039$}; \draw[ink, line width=1.6pt] ([xshift=-0.10cm]eq.west) -- ([xshift=0.10cm]eq.east); \fill[ink] (0.15,0.075) circle (0.130); \draw[white, line width=0.9pt, line cap=round] (0.085,0.010) -- (0.215,0.140); \draw[white, line width=0.9pt, line cap=round] (0.085,0.140) -- (0.215,0.010); }\\[3pt] is the subtraction the shared chance value invites, and it is not a quantity: two different frames, measured against two different references.};
\node[fill=panelbg, inner xsep=3pt, inner ysep=1pt, anchor=west, font=\sffamily\scriptsize\bfseries, text=ink] at ([xshift=0.30cm]nosub.north west) {Do not subtract};

% ================ the scope limit, drawn as a rule and not as a caption
\draw[hardrule] (0.0300,6.0500) -- (14.1700,6.0500);
\node[anchor=north west, font=\sffamily\scriptsize, text=ink, align=flush left, inner sep=0pt, text width=9.90cm] at (4.30,5.9200) {\textbf{Scope limit.} No number in one frame may be set beside a number in the other.};
% the connector is stopped at the rule, with the panel's own mark
\fill[ink] (3.8894,6.0500) circle (0.190);
\draw[white, line width=1.0pt, line cap=round] (3.7944,5.9550) -- (3.9844,6.1450);
\draw[white, line width=1.0pt, line cap=round] (3.7944,6.1450) -- (3.9844,5.9550);

% ================ LOWER SCALE -- residual frame
\node[fh,  anchor=west] at (0.0000,5.05) {residual frame};
\node[fhs, anchor=east] at (14.2000,5.05) {cell-line sets, $n=\num{1394}$};
\draw[headrule] (0.0300,4.75) -- (14.1700,4.75);
\draw[scaleline] (2.35,3.6000) -- (12.30,3.6000);
\fill[quietink] (3.6000,3.6000) circle (1.4pt);   % generic
\fill[quietink] (3.6113,3.6000) circle (1.4pt);   % scramble_orth
\fill[quietink] (3.6126,3.6000) circle (1.4pt);   % random
\fill[quietink] (3.7118,3.6000) circle (1.4pt);   % control_copy
\fill[ink]    (10.6000,3.6000) circle (1.6pt);   % within-well reference
\fill[ink]    (11.7597,3.6000) circle (1.6pt);   % drug_lookup
\fill[accent] (4.3259,3.6000) circle (1.9pt);   % model
\draw[accent!85, line width=0.5pt] (3.6000,3.4900) -- (3.6000,3.7100);
\draw[leader] (4.3259,3.6900) -- (4.46,3.8400);
\node[val, anchor=south west, align=left, text=accent, font=\footnotesize\bfseries] at (4.44,3.7800) {model\\ $0.537$};
\node[val, anchor=south] at (10.6000,3.7500) {$0.854$};
\draw[leader] (11.7597,3.6900) -- (11.98,3.8400);
\node[arm, anchor=south west, align=left, text=ink] at (11.96,3.7800) {\texttt{drug\_lookup}\\ $0.913$};
\draw[leader] (3.6818,3.5100) -- (3.44,3.26);
\draw[leader] (3.40,3.26) -- (3.40,2.29);
\draw[leader] (3.40,3.2600) -- (3.34,3.2600);
\node[arm, anchor=east] at (3.30,3.2600) {\texttt{generic}\;$0.500$};
\draw[leader] (3.40,2.9400) -- (3.34,2.9400);
\node[arm, anchor=east] at (3.30,2.9400) {scramble, \texttt{orth}\;$0.501$};
\draw[leader] (3.40,2.6200) -- (3.34,2.6200);
\node[arm, anchor=east] at (3.30,2.6200) {\texttt{random}\;$0.501$};
\draw[leader] (3.40,2.3000) -- (3.34,2.3000);
\node[arm, anchor=east] at (3.30,2.3000) {\texttt{control\_copy}\;$0.506$};
\draw[leader] (4.3659,3.5100) -- (4.66,3.30);
\node[note, anchor=north west, text width=5.50cm] at (4.70,3.36) {The $13.9\%$ margin this chapter quotes for the model (\cref{fig:ladder}) is computed on the $n=\num{1192}$ common support, where the model is $0.549$ against a reference of $0.857$. It is not a measurement of the $n=\num{1394}$ gap drawn here.};
\draw[calib] (13.1810,3.00) -- (14.1700,3.00);
\draw[calib] (13.1810,2.92) -- (13.1810,3.08);
\draw[calib] (14.1700,2.92) -- (14.1700,3.08);
\node[noteR, anchor=north east] at (14.2000,2.86) {$0.05$ of $\NIR$ on this scale};
\node[noteR, anchor=north east, text width=3.85cm] at (14.2000,2.44) {The two rulers differ by $4.6\times$.};
\end{tikzpicture}
%
% ---------------------------------------------------------------------------
% READY TO PASTE into Sections/04d-baselines.tex, at sec:res-summary, just after
% the paragraph "The same chance value, two different references".
%
% \begin{figure}[htbp]
% \centering
% \input{figs/two-scales}
% \caption[Two frames, one chance value]{\textbf{The two frames share a chance
% value and nothing else, and the shared value is exactly what makes the mistake
% easy.} Each scale is stretched so that chance falls at the same place in both
% and so that each frame's own within-well split-half precision reference falls
% at the same place as the other's. Nothing else is aligned, so the two frames
% carry different rulers --- the same $0.05$ of $\NIR$ is 4.6 times longer
% above than below --- and that differential stretch is the figure.
% \emph{Upper:} the expression frame, within plate, tier 2, $n = \num{606}$,
% where the model scores $0.498$ against a reference of $0.576$ and is matched by
% a drug-agnostic linear map ($0.500$) and a control-copy ($0.504$); the scale is
% broken at the left because the global mean, $0.180$, lies $4.2\times$ that
% frame's chance-to-reference distance below chance. \emph{Lower:} the residual
% frame, cell-line sets, $n = \num{1394}$, where the model scores $0.537$ against
% a reference of $0.854$, a training-only per-drug lookup reaches $0.913$ above
% that reference, and \texttt{generic}, \texttt{orth}, \texttt{random} and
% \texttt{control\_copy} all fall within $0.006$ of chance, one mark at this
% scale. Both references are within-well split-half precision figures --- two
% disjoint halves of one treated well --- and not biological replicates
% (\cref{sec:methods-data}). Chance is $0.50$ in both under the exchangeability
% condition audited in \cref{sec:res-metrics}, and because it is the one value
% the two frames share it is exactly what makes the struck-out subtraction in the
% middle of the figure easy to perform: $0.537$ and $0.498$ are not two attempts
% at one task, and the difference between them is neither an improvement nor a
% decline. Every value is a pooled mean drawn as a bare marker; no arm shown here
% is quoted with an interval at the pooled level, so none is drawn
% (\cref{tab:allarms}). The $13.9\%$ margin annotated below is computed on the
% $n = \num{1192}$ common support and is not a measurement of the
% $n = \num{1394}$ gap drawn here.}
% \label{fig:two-scales}
% \end{figure}
% ---------------------------------------------------------------------------

% \caption[Two frames, one chance value]{\textbf{The two frames share a chance
% value and nothing else, and the shared value is exactly what makes the mistake
% easy.} Each scale is stretched so that chance falls at the same place in both
% and so that each frame's own within-well split-half precision reference falls
% at the same place as the other's. Nothing else is aligned, so the two frames
% carry different rulers --- the same $0.05$ of $\NIR$ is 4.6 times longer
% above than below --- and that differential stretch is the figure.
% \emph{Upper:} the expression frame, within plate, tier 2, $n = \num{606}$,
% where the model scores $0.498$ against a reference of $0.576$ and is matched by
% a drug-agnostic linear map ($0.500$) and a control-copy ($0.504$); the scale is
% broken at the left because the global mean, $0.180$, lies $4.2\times$ that
% frame's chance-to-reference distance below chance. \emph{Lower:} the residual
% frame, cell-line sets, $n = \num{1394}$, where the model scores $0.537$ against
% a reference of $0.854$, a training-only per-drug lookup reaches $0.913$ above
% that reference, and \texttt{generic}, \texttt{orth}, \texttt{random} and
% \texttt{control\_copy} all fall within $0.006$ of chance, one mark at this
% scale. Both references are within-well split-half precision figures --- two
% disjoint halves of one treated well --- and not biological replicates
% (\cref{sec:methods-data}). Chance is $0.50$ in both under the exchangeability
% condition audited in \cref{sec:res-metrics}, and because it is the one value
% the two frames share it is exactly what makes the struck-out subtraction in the
% middle of the figure easy to perform: $0.537$ and $0.498$ are not two attempts
% at one task, and the difference between them is neither an improvement nor a
% decline. Every value is a pooled mean drawn as a bare marker; no arm shown here
% is quoted with an interval at the pooled level, so none is drawn
% (\cref{tab:allarms}). The $13.9\%$ margin annotated below is computed on the
% $n = \num{1192}$ common support and is not a measurement of the
% $n = \num{1394}$ gap drawn here.}
% \label{fig:two-scales}
% \end{figure}
% ---------------------------------------------------------------------------

% \caption[Two frames, one chance value]{\textbf{The two frames share a chance
% value and nothing else, and the shared value is exactly what makes the mistake
% easy.} Each scale is stretched so that chance falls at the same place in both
% and so that each frame's own within-well split-half precision reference falls
% at the same place as the other's. Nothing else is aligned, so the two frames
% carry different rulers --- the same $0.05$ of $\NIR$ is 4.6 times longer
% above than below --- and that differential stretch is the figure.
% \emph{Upper:} the expression frame, within plate, tier 2, $n = \num{606}$,
% where the model scores $0.498$ against a reference of $0.576$ and is matched by
% a drug-agnostic linear map ($0.500$) and a control-copy ($0.504$); the scale is
% broken at the left because the global mean, $0.180$, lies $4.2\times$ that
% frame's chance-to-reference distance below chance. \emph{Lower:} the residual
% frame, cell-line sets, $n = \num{1394}$, where the model scores $0.537$ against
% a reference of $0.854$, a training-only per-drug lookup reaches $0.913$ above
% that reference, and \texttt{generic}, \texttt{orth}, \texttt{random} and
% \texttt{control\_copy} all fall within $0.006$ of chance, one mark at this
% scale. Both references are within-well split-half precision figures --- two
% disjoint halves of one treated well --- and not biological replicates
% (\cref{sec:methods-data}). Chance is $0.50$ in both under the exchangeability
% condition audited in \cref{sec:res-metrics}, and because it is the one value
% the two frames share it is exactly what makes the struck-out subtraction in the
% middle of the figure easy to perform: $0.537$ and $0.498$ are not two attempts
% at one task, and the difference between them is neither an improvement nor a
% decline. Every value is a pooled mean drawn as a bare marker; no arm shown here
% is quoted with an interval at the pooled level, so none is drawn
% (\cref{tab:allarms}). The $13.9\%$ margin annotated below is computed on the
% $n = \num{1192}$ common support and is not a measurement of the
% $n = \num{1394}$ gap drawn here.}
% \label{fig:two-scales}
% \end{figure}
% ---------------------------------------------------------------------------

\caption[Two frames, one chance value]{\textbf{The two frames share a chance
value and nothing else, and the shared value is exactly what makes the mistake
easy.} Each scale is stretched so that chance falls at the same place in both and
so that each frame's own within-well split-half precision reference falls at the
same place as the other's; nothing else is aligned, so the two frames carry
different rulers --- the same $0.05$ of $\NIR$ is $4.6\times$ longer above than
below --- and that differential stretch is the figure. \emph{Upper:} the
expression frame, within plate, tier 2, $n = \num{606}$, where the model scores
$0.498$ against a reference of $0.576$ and is matched by a drug-agnostic linear
map ($0.500$) and a control-copy ($0.504$); the scale is broken at the left
because the global mean, $0.180$, lies $4.2\times$ that frame's
chance-to-reference distance below chance. \emph{Lower:} the residual frame,
cell-line sets, $n = \num{1394}$, where the model scores $0.537$ against a
reference of $0.854$, a training-only per-drug lookup reaches $0.913$ above that
reference, and \texttt{generic}, \texttt{orth}, \texttt{random} and
\texttt{control\_copy} all fall within $0.006$ of chance, one mark at this scale.
Both references are within-well split-half precision figures --- two disjoint
halves of one treated well --- and not biological replicates
(\cref{sec:methods-data}). Chance is $0.50$ in both under the exchangeability
condition audited in \cref{sec:res-metrics}, and because it is the one value the
two frames share it is exactly what makes the struck-out subtraction in the
middle of the figure easy to perform. Every value is a pooled mean drawn as a
bare marker; no arm shown here is quoted with an interval at the pooled level, so
none is drawn (\cref{tab:allarms}). The $13.9\%$ margin annotated below is
computed on the $n = \num{1192}$ common support and is not a measurement of the
$n = \num{1394}$ gap drawn here.}
\label{fig:two-scales}
\end{figure}

% PLACEMENT ONLY. The ledger was anchored two paragraphs and a scope-limit box
% lower down, past the point where a float 0.82 of the column high can still be
% taken: it was deferred, took a page to itself, and left the three lines that
% close the chapter alone on the page after it. Anchored here, at the sentence
% that announces it -- "What follows is a ledger of the principal comparison
% arms" -- it is offered a page earlier and heads a text page instead. Reading
% order is unchanged: the table still sets before the prose that reads it.
% Chapter 4 previously never \cref'd this table -- the nearest references were in
% ch.5 and the appendix, five and thirty-odd pages away. It is now \cref'd on
% both sides of the float, one page either side: in the frame-defining paragraph
% immediately above (the page before the table sets) and in "What the ledger
% deliberately leaves out" below. Verified page-neutral against a pristine build:
% same section page, same table page, same chapter-end page, same page count.
% ---------------------------------------------------------------------------
\begin{table}[htbp]
\begin{fullwidth}
\begin{threeparttable}
\caption{\textbf{The principal comparison arms on the calibrated metric, by
frame.} Chance is $0.50$ in both; reference rows are within-well split-half
precision figures, not biological replicates (\cref{sec:methods-data}). The final
expression rows show the $204$ conditions clearing the pre-set $0.80$
identifiability threshold, where a drug-agnostic control still matches the model.
Of the scramble strata, \texttt{near} ($0.550$) and \texttt{opposite} ($0.423$)
differ from chance in opposite directions while \texttt{orth} ($0.501$) is
empirically near chance here; because the strata are defined using true residual
geometry, \texttt{orth} is the least contaminated available comparator and not an
independently established null (\cref{sec:res-generalisation}).}
\label{tab:allarms}
\small
\begin{tabular}{llr>{\raggedright\arraybackslash}p{0.30\linewidth}}
\toprule
\thead{Frame} & \thead{Arm} & {\thead{$\NIR$}} & \thead{What it knows} \\
\midrule
\multirow{7}{*}{\shortstack[l]{expression\\(within plate,\\tier 2, $n=606$)}}
 & within-well reference & 0.576 & two disjoint halves of one treated well;
precision, not a replicate \\
 & control-copy & 0.504 & nothing; it reproduces the untreated cell \\
 & linear map & 0.500 & the control cell, no drug \\
 & \textbf{model} & \textbf{0.498} & control cell $+$ drug \\
 & global mean & 0.180 & nothing \\
 & \textbf{model}, identifiable stratum & \textbf{0.768} & the same model, on the
$204$ conditions whose own reference clears $0.80$ \\
 & control-copy, same stratum & 0.766 & nothing, and it matches the model
there \\
\midrule
\multirow{7}{*}{\shortstack[l]{residual\\(cell-line sets,\\$n=1394$)}}
 & within-well reference & 0.854 & two disjoint halves of one treated well; not a
replicate, and on the lookups' support, where it is $0.857$, the lookup exceeds
it \\
 & drug lookup (training-only) & 0.913 & this drug's residual from other cell
lines, on the $n = \num{1192}$ conditions admitting a lookup, where the model
scores $0.549$ \\
 & drug lookup, one line & 0.822 & the same, from one other cell line, same
support \\
 & \textbf{model} & \textbf{0.537} & control cell $+$ drug \\
 & scramble, \texttt{orth} partner & 0.501 & geometry-defined by residual cosine
to the target ($\cos \approx 0$); its own score is \emph{empirically} near chance
here \\
 & scramble, \texttt{opposite} partner & 0.423 & an active treatment, not a
null \\
 & generic response & 0.500 & nothing (the zero vector, by construction) \\
\bottomrule
\end{tabular}
\end{threeparttable}
\end{fullwidth}
\end{table}

\paragraph{The same chance value, two different references} Both use $0.50$ as
chance under the exchangeability condition audited in \cref{sec:res-metrics}.
\Cref{fig:two-scales} puts each frame on its own ruler, scaled so that chance and
that frame's own within-well precision reference fall at the same two places: the
$0.537$ in one frame and
the $0.498$ in the other are not two attempts at one task, and the difference
between them is neither an improvement nor a decline. Read each against its own
reference, the expression-frame model does not clear chance at all, and the
residual-frame model clears it on trained conditions only,
\ci{0.598}{0.549}{0.647} there against \ci{0.533}{0.492}{0.575} and
\ci{0.459}{0.404}{0.514} on the held-out splits, neither distinguishable from
$0.50$ (\cref{sec:res-generalisation}). The pooled $0.537$ is three-quarters
held-out condition, and no arm is quoted with an interval at the pooled level, so
it is quoted as a mean and not as a clearance.


\paragraph{And the margin between them has a support} What separates the frames
is a margin \cref{sec:res-lookup} sizes at $13.9\%$ of the distance from chance
to the reference, computed on the $n = \num{1192}$ common support where the model
scores $0.549$ against a reference of $0.857$, and not on the $n = \num{1394}$
whose $0.537$ and $0.854$ are quoted above.

\begin{scopelimit}[carried into \cref{sec:limitations}]
No number in one frame may be set beside a number in the other. Chance is $0.50$
in both, which is exactly what makes the mistake easy to make; the references
above chance are $0.854$ and $0.576$, and every distance quoted in this chapter is
a distance to its own frame's reference. Within the residual frame the same
applies to supports: $\num{1394}$, $\num{1192}$ and $\num{218}$ select conditions
by different rules, and rows on different supports are not comparable.
\end{scopelimit}

% ---------------------------------------------------------------------------
\paragraph{What the ledger deliberately leaves out} \cref{tab:allarms} omits the
\texttt{near} scramble, a deliberately weak test; \texttt{moa\_lookup}, whose
$n=218$ support is not comparable with the principal rows; the residual-frame
\texttt{control\_copy} and \texttt{random} negatives, which serve the calibration
diagnostic; and the field-attribution arms, whose non-neutral comparators leave
attribution unresolved. The expression-frame control-copy stays because it
carries the identifiable-stratum comparison, and \texttt{orth} and
\texttt{opposite} because their contrast exposes how comparator score depends on
partner geometry.

\paragraph{Five results are estimates, not predictor arms} They are stated
outside the ledger because nothing in it is their comparison set.

\begin{itemize}
\item \textbf{The metric audit.} Of the five metrics graded, $\NIR$ is the only
one whose $\DRF$ is positive, \ci{+0.635}{+0.604}{+0.663} across plates,
\ci{+0.545}{+0.516}{+0.569} within plate, \ci{+0.704}{+0.667}{+0.736} in the
residual frame; the other four are negative and exclude zero. Three,
$\paneltau$, \texttt{spearman\_expr} and $\DEdr$, are inverted rather than merely
uninformative, scoring an uninformative baseline above the positive control;
\texttt{weighted\_r2} sits closest to zero and its sign tracks cell count, so it
carries no directional claim (\cref{sec:res-metrics}).
\item \textbf{Representation and readout.} Drug identity is decodable from the
activations, and the generation readout consults it at one to two per cent of the
preference gap, a fifth or less of the pre-registered material margin, in one
training arm and at one rung of the displacement ladder; the chapter's weakest
load-bearing result (\cref{sec:res-mechanism}).
\item \textbf{Prompt sensitivity.} Re-encoding the target as the drug-specific
residual gives \ci{+0.0898}{+0.0291}{+0.1506} on trained conditions, clustered on
the drug. A budget-matched optimal-transport control clears zero too, at
\ci{+0.0292}{+0.0144}{+0.0441}, so the re-encoding is first in the sequence and
not uniquely responsible (\cref{sec:res-encoding}).
\item \textbf{Transfer.} \ci{\Tcoef = 0.553}{0.514}{0.593}, a cross-context
similarity coefficient and not a variance share (\cref{sec:res-transfer}).
\item \textbf{The channels.} Same-target retrieval is the strongest observed lead
on its restricted support; only that arm clears the relevance margin, on the
cell-line axis on all three constructions and not on the drug axis under either
estimand-correct null, on $18$ annotated drugs of which two are held out. The
arms differ in support, roster, partner count and annotation coverage, so this is
not an identified biological ranking (\cref{sec:res-scope}).
\end{itemize}

\noindent
The distance between the model and the least contaminated available scramble is
the prompt-sensitivity effect this work recovered; the distance between the model
and the lookup is what it did not.
